% concatenated from main.tex

\documentclass[12pt, reqno]{amsart}
\usepackage{amsmath, amsthm, amscd, amsfonts, amssymb, graphicx, color}
\usepackage[bookmarksnumbered, colorlinks, plainpages]{hyperref}
\allowdisplaybreaks

% 页面设置
\textheight 22.5truecm \textwidth 14.5truecm
\setlength{\oddsidemargin}{0.35in}\setlength{\evensidemargin}{0.35in}
\setlength{\topmargin}{-.5cm}
\usepackage{microtype} 



\hypersetup{
    colorlinks=true,    % 使用彩色链接（而非边框）
    linkcolor=black,    % 目录和引用链接改为黑色
    filecolor=black,
    urlcolor=blue,      % 网址保持蓝色或改为黑色
    citecolor=black,    % 引用编号改为黑色
}

\usepackage{titletoc}

\titlecontents{section}
  [0em]                              % 左缩进
  {\vspace{0.5em}}                   % 间距
  {\thecontentslabel\hspace{0.8em}}  % 编号 + 间距（关键！）
  {}                                 % 无编号条目
  {\titlerule*[0.5pc]{.}\contentspage} % 引导线 + 页码

\titlecontents{subsection}
  [1.5em]
  {}
  {\thecontentslabel\hspace{0.8em}}  % 同样添加 \hspace
  {}
  {\titlerule*[0.5pc]{.}\contentspage}

\newtheorem{theorem}{Theorem}[section]
\newtheorem{lemma}[theorem]{Lemma}
\newtheorem{proposition}[theorem]{Proposition}
\newtheorem{corollary}[theorem]{Corollary}
\theoremstyle{definition}
\newtheorem{definition}[theorem]{Definition}
\newtheorem{example}[theorem]{Example}
\newtheorem{exercise}[theorem]{Exercise}
\newtheorem{conclusion}[theorem]{Conclusion}
\newtheorem{conjecture}[theorem]{Conjecture}
\newtheorem{criterion}[theorem]{Criterion}
\newtheorem{summary}[theorem]{Summary}
\newtheorem{axiom}[theorem]{Axiom}
\newtheorem{problem}[theorem]{Problem}
\theoremstyle{remark}
\newtheorem{remark}[theorem]{Remark}
\numberwithin{equation}{section}
\usepackage{titlesec}
\titleformat{\section}[hang]{\normalfont\Large\bfseries}{\thesection.}{0.5em}{}
\titleformat{\subsection}[hang]{\normalfont\large\bfseries}{\thesubsection.}{0.5em}{}
\titleformat{\subsubsection}[hang]
  {\normalfont\normalsize\bfseries}
  {\thesubsubsection.}
  {0.5em}
  {}
\DeclareMathOperator{\Irred}{Irred}
\DeclareMathOperator{\Pol}{Pol}
  
\begin{document}
\setcounter{page}{1}



\centerline{}

\centerline{}



%Title of the paper
\title[GH convergence of spectral truncation for  compact quantum groupS]{Gromov-Hausdorff Convergence of Spectral Truncations for Compact  Quantum GroupS
}

%Author names and affiliations
\author{Xintao Peng , Qin Wang}

%In case of 3 or more authors use below format
%\author[F. Author, S. Author, T. Author]{First Author$^1$, Second Author$^2$$^{*}$ \MakeLowercase {and} Third Author$^3$}



%Abstract, keywords, math subject classification

\begin{abstract}


We study the quantum Gromov–Hausdorff convergence of spectral truncations for compact quantum groups. Using a proper length function, we define a Dirac operator and the associated spectral truncations. This work extends the previous convergence results for tori (Leimbach–van Suijlekom) to a broad class of quantum groups, and provides Gromov–Hausdorff convergence result for spectral truncations on quantum groups, encompassing both compact  and discrete quantum groups. Our results are applicable to $SU(N)$,$SO(N)$ and discrete quantum groups of rapid decay.
\end{abstract}

\maketitle

%------------------------------------------------------------------------------

\section{Introduction}







Rieffel~\cite{9} introduced compact quantum metric spaces, extending the Gromov--Hausdorff distance to the noncommutative setting. Since then, spectral truncations in noncommutative geometry have been studied by Connes and van Suijlekom~\cite{1}, later by van Suijlekom~\cite{5}, and by Leimbach and van Suijlekom~\cite{4}, who established convergence results for noncommutative tori. Rieffel~\cite{11} also considered Fourier truncations for compact quantum groups and finitely generated groups, while Leimbach~\cite{6} studied the convergence of Peter--Weyl truncations for compact quantum groups. Nevertheless, a systematic treatment of spectral truncations for compact quantum groups has remained largely unavailable.

A crucial step toward such a theory was made by Austad and Kyed~\cite{16}. Let $\mathbb{G}$ be a compact quantum group equipped with a proper length function $\ell:\Irred(\mathbb{G})\to[0,\infty)$. As shown in~\cite{23}, one obtains a spectral triple
\[
\bigl(\Pol(\mathbb{G}),\, L^{2}(\mathbb{G}),\, D_{\ell}\bigr),
\]
where $\Pol(\mathbb{G})$ is the Hopf $*$-algebra of representative functions, $L^{2}(\mathbb{G})$ is the GNS Hilbert space of the Haar state, and $D_{\ell}$ acts on matrix coefficients by
\[
D_{\ell}\bigl(\Lambda(u_{i,j}^{\alpha})\bigr)=\ell(\alpha)\,\Lambda(u_{i,j}^{\alpha}).
\]
However, this spectral triple does not automatically define a compact quantum metric space; one must additionally verify that the associated Lipschitz seminorm induces the weak$^{*}$ topology on the state space.

Austad and Kyed proved that, for coamenable compact quantum groups of Kac type, compact quantum metric properties of the full algebra are equivalent to those of the central subalgebra. More precisely, if $L_{\ell}$ denotes the Lip-norm associated with $\ell$ and $L_{z,\ell}$ its restriction to the central part, then $(C(\mathbb{G}),L_{\ell})$ is a compact quantum metric space if and only if $(C_{z}(\mathbb{G}),L_{z,\ell})$ is. This reduction allows one to study the metric structure through central functions or, equivalently, through the fusion algebra
\[
F(\mathbb{G})=\operatorname{span}_{\mathbb{C}}\{\alpha:\alpha\in\Irred(\mathbb{G})\}.
\]
In particular, all three settings---$C(\mathbb{G})$, $C_{z}(\mathbb{G})$, and $F(\mathbb{G})$---carry compatible metric information.

The present paper studies spectral truncations for compact quantum groups and their convergence in quantum Gromov--Hausdorff distance. We first recall the definition of quantum Gromov--Hausdorff distance between compact quantum metric spaces.

\begin{definition}[Quantum Gromov--Hausdorff distance]
Let $(A,L_{A})$ and $(B,L_{B})$ be compact quantum metric spaces. Let $\mathcal{M}(L_{A},L_{B})$ denote the set of Lip-norms on $A\oplus B$ that induce $L_{A}$ and $L_{B}$. The quantum Gromov--Hausdorff distance between $A$ and $B$ is
\[
\operatorname{dist}_{q}(A,B)
=
\inf\Bigl\{
\operatorname{dist}_{H}^{d^{L}}\bigl(S(A),S(B)\bigr)
:\, L\in\mathcal{M}(L_{A},L_{B})
\Bigr\}.
\]
\end{definition}

We then consider spectral truncations associated with a spectral triple $(A,\mathcal{H},D)$, where $A$ acts on a Hilbert space $\mathcal{H}$ and $D$ is self-adjoint with compact resolvent. For a cutoff $\Lambda$, let $P_{\Lambda}$ be the spectral projection of $D$ onto the eigenspaces with eigenvalues of modulus at most $\Lambda$. The truncated algebra is $P_{\Lambda}AP_{\Lambda}$ acting on $P_{\Lambda}\mathcal{H}$, with truncated operator $D_{\Lambda}=D|_{P_{\Lambda}\mathcal{H}}$. Our goal is to determine when these truncations converge to the original algebra in quantum Gromov--Hausdorff distance.

For compact quantum groups, we focus on the fusion algebra
\[
F(\mathbb{G})=\operatorname{span}_{\mathbb{C}}\{\alpha:\alpha\in\Irred(\mathbb{G})\},
\]
equipped with the natural $\ell^{2}$-structure coming from the dimensions $d(\alpha)$. We say that $F(\mathbb{G})$ has the rapid decay property if there exist constants $C>0$ and $s>0$ such that
\[
\|\Lambda(f)\|\le C\|f\|_{2,s}
\]
for all $f\in F(\mathbb{G})$. We also introduce polynomial growth and strong polynomial growth with respect to a length function $\ell$.

Our first main result shows that polynomial growth yields compact quantum metric structures on higher commutator seminorms.

\begin{theorem}
If $\mathbb{G}$ is a compact quantum group with polynomial growth of order $s$, then for every integer $k>s$, the pair $\bigl(F(\mathbb{G}),L_{\ell}^{k}\bigr)$ is a compact quantum metric space.
\end{theorem}

This result applies, in particular, to connected compact Lie groups such as $SU(N)$ and $SO(N)$ equipped with suitable Lip-norms.

We next study spectral truncations of $C(\mathbb{G})$. Let
\[
C(\mathbb{G})_\Lambda=P_{\Lambda}C(\mathbb{G})P_{\Lambda},
\]
and define the seminorms $L_{C(\mathbb{G})_\Lambda}^{\alpha^\tau}$, $L_{C(\mathbb{G})_\Lambda}^{\beta^\tau}$, and $L_{C(\mathbb{G})_\Lambda}^{\alpha^\tau,\beta^\tau}$ from the left and right coactions as above. Our convergence theorem is as follows.

\begin{theorem}\label{1}
Let $\mathbb{G}$ be a compact quantum group equipped with a proper length function. If $\bigl(C(\mathbb{G}),L_{\ell}^{k}\bigr)$ is a compact quantum metric space, then 
\[
\operatorname{dist}_{q}\Bigl(
\bigl(C(\mathbb{G}),L_{\ell}^{k}\bigr),\,
\bigl(P_{\Lambda}C(\mathbb{G})P_{\Lambda},L_{C(\mathbb{G})_\Lambda}^{\alpha^\tau,\beta^\tau}\bigr)
\Bigr)\to 0 \qquad as \quad\Lambda\to \infty.
\]
\end{theorem}

The proof constructs suitable maps $\tau:A\to A_{\Lambda}$ and $\sigma^{\phi}:A_{\Lambda}\to A$ and verifies the hypotheses of Definition~2 in~\cite{5}.

We then turn to the central subalgebra
\[
C_{z}(\mathbb{G})
=
\{a\in C(\mathbb{G})\mid \sigma\Delta(a)=\Delta(a)\},
\]
where $\sigma$ is the flip map on $C(\mathbb{G})\otimes C(\mathbb{G})$. Extending the result of Austad and Kyed~\cite{16} from $k=1$ to arbitrary $k\in\mathbb{N}$, we prove that for a compact, coamenable quantum group of Kac type and a proper length function $\ell$, then $(C(\mathbb{G}), L_{\ell}^{k})$ is a compact quantum metric space if and only if $(C_{z}(\mathbb{G}), L_{z,\ell}^{k})$ is a compact quantum metric space.

Finally, we apply the same method to discrete quantum groups of rapid decay. In this setting, we obtain analogous quantum Gromov--Hausdorff convergence for the corresponding spectral truncations of $C(\widehat{\mathbb{G}})$.

The paper is organized as follows. Section~2 recalls the necessary background on quantum Gromov--Hausdorff convergence, compact quantum groups, and fusion algebras. Section~3 studies compact quantum metric spaces arising from fusion algebras with rapid decay. Section~4 proves quantum Gromov--Hausdorff convergence for truncations of $C(\mathbb{G})$ under the assumption that $(C(\mathbb{G}),L_{\ell}^{k})$ is a compact quantum metric space. Section~5 proves for a coamenable Kac-type quantum group $\mathbb{G}$, the pair
$\bigl(C(\mathbb{G}), L_{\ell}^{k}\bigr)$ is a compact quantum metric space
if and only if $\bigl(C_{z}(\mathbb{G}), L_{z,\ell}^{k}\bigr)$ is a compact
quantum metric space,
 and also treats the discrete quantum group case.













\section{Preliminaries}

\subsection{Quantum Gromov-Hausdorff Distance}
This subsection reviews the concept of compact quantum metric spaces 
and the quantum analogue of Gromov--Hausdorff convergence introduced 
by Marc A.~Rieffel in \cite{9}. The relevant definitions and key 
properties are recalled below.

\begin{definition}\label{Definition2.1}
Let $A$ be a complete operator system. A seminorm 
$L \colon A \to [0,\infty]$ is called a \emph{Lipschitz seminorm} 
if $L$ satisfies the following conditions:
\begin{itemize}
    \item[(1)] $L$ is \emph{densely defined}, meaning that the domain 
    $\operatorname{Dom}(L) := \{a \in A : L(a) < \infty\}$ is a 
    norm-dense subspace of $A$;
    
    \item[(2)] $L(a^*) = L(a)$ for all $a \in A$, and the kernel of $L$ 
    is exactly the scalar multiples of the identity, i.e.\ 
    $L(a) = 0 \iff a \in \mathbb{C}\cdot e_A$.
\end{itemize}
\end{definition}

\begin{definition}[Compact Quantum Metric Space]\label{Definition2.2}
Let $A$ be an operator system equipped with a Lipschitz seminorm 
$L \colon A \to [0,+\infty]$. We define a metric $d^L$ on the state 
space $\mathcal{S}(A)$ by
\begin{equation}\label{eq:dist}\tag{1}
    d^L(\phi, \psi) 
    := \sup\bigl\{ |\phi(a) - \psi(a)| : a \in A,\ L(a) \leq 1 \bigr\},
    \qquad \phi, \psi \in \mathcal{S}(A).
\end{equation}
We require that the topology induced by $d^L$ on $\mathcal{S}(A)$ coincides with 
the weak$^{*}$ topology. The pair $(A, L)$ is then called a 
\emph{compact quantum metric space}. In this case, $L$ is referred to 
as a \emph{Lip-norm}.
\end{definition}



\begin{remark}
The metric $d^L$ defined in (1) is referred to as the 
Monge--Kantorovi\v{c} metric. The condition that $d^L$ metrises 
the weak$^{*}$ topology ensures that $\mathcal{S}(A)$ is compact, hence the 
terminology.
\end{remark}

Next, we introduce the \emph{quantum Gromov--Hausdorff distance} 
between two compact quantum metric spaces.

\begin{definition}\label{Definition2.4}
Let $(A, L_A)$ and $(B, L_B)$ be compact quantum metric spaces. 
Let $\mathcal{M}(L_A, L_B)$ be the set of Lip-norms on $A \oplus B$ 
which induce $L_A$ and $L_B$. The \emph{quantum Gromov--Hausdorff 
distance} between them, denoted $\operatorname{dist}_q(A, B)$, is 
defined by
\[
\operatorname{dist}_q(A, B) = \inf\bigl\{ 
\operatorname{dist}_H^{d^L}(\mathcal{S}(A), \mathcal{S}(B)) : L \in \mathcal{M}(L_A, L_B) 
\bigr\}.
\]
\end{definition}

We now present two propositions that provide an upper bound for the 
quantum Gromov--Hausdorff distance.

\begin{proposition}[\cite{9}, Proposition~8.5]\label{prop:upper_bound}
Let $(A, L_A)$ be a compact quantum metric space, and let $B$ be a 
subspace of $A$ containing $e_A$. Let $L_B$ be the restriction of 
$L_A$ to $B$, so that $(B, L_B)$ is a compact quantum metric space. 
Let $P : A \to B$ be a map 
for which there exists a $\delta > 0$ such that
\begin{enumerate}
    \item $L_B(P(a)) \leq L_A(a)$ for all $a \in A$;
    \item $\|a - P(a)\| \leq \delta\, L_A(a)$ for all $a \in A$.
\end{enumerate}
Then $\operatorname{dist}_q(A, B) \leq \delta$.
\end{proposition}

\begin{definition}
    Let $(A, L_A)$ be a compact quantum metric space with state space $\mathcal{S}(A)$. The \emph{diameter} is given by
    \[
        \operatorname{diam}(A, L_A) 
        := \sup \Bigl\{ d^{L_A}(\mu, \nu) \;\big|\; \mu, \nu \in \mathcal{S}(A) \Bigr\}.
    \]
\end{definition}
\begin{proposition}\label{prop:gromov hausdorff converges}
 
Let $(A, L_A)$ and $(B, L_B)$ be compact quantum metric spaces and suppose that $\Phi: A \to B$ and $\Psi: B \to A$ are two unital positive maps satisfying that
\begin{enumerate}
    \item[(1)] there exist $C, C' > 0$ such that
    \[
    L_B(\Phi(a)) \leqslant C \cdot L_A(a) \quad \text{and} \quad L_A(\Psi(b)) \leqslant C' \cdot L_B(b)
    \]
    for all $a \in A$ and $b \in B$;
    \item[(2)] there exist $\varepsilon, \varepsilon' > 0$ such that
    \[
    \|\Psi\Phi(a) - a\| \leqslant \varepsilon \cdot L_A(a) \quad \text{and} \quad \|\Phi\Psi(b) - b\| \leqslant \varepsilon' \cdot L_B(b)
    \]
    for all $a \in A$ and $b \in B$.
\end{enumerate}
Then the quantum Gromov--Hausdorff distance $\operatorname{dist}_{q}(A,B)$ is dominated by
\[
\max\bigl\{ \operatorname{diam}(A, L_A) \cdot |1 - 1/C| + \varepsilon/C, \operatorname{diam}(B, L_B) \cdot |1 - 1/C'| + \varepsilon'/C' \bigr\}.
\]
\end{proposition}

\begin{proof}
To ease the notation, we put
\[
r= \max\bigl\{ \operatorname{diam}(A, L_A) \cdot |1 - 1/C| + \varepsilon/C, \operatorname{diam}(B, L_B) \cdot |1 - 1/C'| + \varepsilon'/C' \bigr\},
\]
and define a Lipschitz seminorm $L: A \oplus B \to [0, \infty]$ by
\[
L(a, b) = \max\left\{ L_A(a), L_B(b), \frac{1}{r}\|b - \Phi(a)\|, \frac{1}{r}\|a - \Psi(b)\| \right\}.
\]

Since both $(A, L_A)$ and $(B, L_B)$ are compact quantum metric spaces, $L$ turns $A \oplus B$ into a compact quantum metric space.

 Let $a \in \operatorname{Dom}(L_A)_{\mathrm{sa}}$ and $\mu \in \mathcal{S}(A)$. Put $z := a - \mu(a)1_A$ and define $b := \frac{1}{C}\Phi(z) + \mu(a)1_B \in \operatorname{Dom}(L_B)_{\mathrm{sa}}$. We estimate:
\begin{enumerate}
    \item $L_B(b) \leqslant \tfrac{1}{C}L_B(\Phi(z)) \leqslant L_A(a)$
    \item $\begin{aligned}[t]
    \tfrac{1}{r}\|b - \Phi(a)\| &= \tfrac{1}{r}\bigl\|\tfrac{1}{C}\Phi(z) - \Phi(z)\bigr\| \\
                               &\leqslant \tfrac{|1-1/C|}{r}\|z\| \\
                               &\leqslant \tfrac{|1-1/C|\operatorname{diam}(A, L_A)}{r} \cdot L_A(a) \\
                               &\leqslant L_A(a)\end{aligned}$
    \item $\begin{aligned}[t]
          \tfrac{1}{r}\|a - \Psi(b)\| &\leqslant \tfrac{1}{r}\bigl\|z - \tfrac{1}{C}\Psi\Phi(z)\bigr\| \\
                               &\leqslant \tfrac{|1-1/C|}{r}\|z\| + \tfrac{1}{Cr}\|z - \Psi\Phi(z)\| \\
                               &\leqslant \tfrac{|1-1/C|\operatorname{diam}(A, L_A)}{r} L_A(a) + \tfrac{1}{C} \cdot \tfrac{\varepsilon}{r} L_A(a)\\ 
                               &\leqslant L_A(a).
          \end{aligned}$
\end{enumerate}
Thus $L(a, b) \leqslant L_A(a)$. The symmetric case for $L_B$ follows similarly.

Finally, for any $\mu \in \mathcal{S}(A)$, put $\nu := \mu \circ \Psi \in  \mathcal{S}(B)$. Then $d^L(\mu, \nu) \leqslant r$. By symmetry,
\[
\operatorname{dist}_{q}(A,B) \leqslant \operatorname{dist}_{\mathrm{H}}^{d^L}( \mathcal{S}(A),  \mathcal{S}(B)) \leqslant r,
\]
which completes the proof.

\end{proof}

Now we focus on spectral triples and the compact quantum metric spaces 
arising from them.

\begin{definition}[Operator system spectral triple]\label{def:op-sys-spectral-triple}
An \emph{operator system spectral triple} is a triple 
$(\mathcal{E}, \mathcal{H}, D)$ where $\mathcal{E}$ is a dense subspace 
of a (concrete) operator system $E$ in $\mathcal{B}(\mathcal{H})$, 
$\mathcal{H}$ is a Hilbert space, and $D$ is a self-adjoint operator 
in $\mathcal{H}$ with compact resolvent such that $[D, T]$ is a bounded 
operator for all $T \in \mathcal{E}$.
\end{definition}

We now consider sequences of operator system spectral triples and their 
limits via spectral truncation.  The limit structure is captured by the following notion of an approximate order isomorphism. 

Let $(\mathcal{E}, \mathcal{H}, D)$ be a spectral triple as above: 
$\mathcal{E}$ acts on $\mathcal{H}$, $D$ is self-adjoint with compact 
resolvent, and $[D, a]$ extends to a bounded operator for all $a$ in a 
dense $*$-subalgebra of $\mathcal{E}$.  
A spectral truncation is given by a spectral projection $P_\Lambda$ of 
$D$ onto the eigenspaces with eigenvalues of modulus $\le \Lambda$.  
The space $\mathcal{E}$ is then truncated to 
$\mathcal{E}_\Lambda := P_\Lambda \mathcal{E} P_\Lambda$, acting on 
$\mathcal{H}_\Lambda := P_\Lambda \mathcal{H}$, and $D$ restricts to the 
self-adjoint operator $D_\Lambda := P_\Lambda D P_\Lambda$ on 
$P_\Lambda \mathcal{H}$.

We define the associated Lipschitz seminorms: for $k \in \mathbb{N}$, 
$a \in \mathcal{E}$ and $b \in \mathcal{E}_\Lambda$, set
\begin{align*}
   & L^k(a)   := \|\delta^k(a)\|   = \bigl\|[D, [D, \dots [D, a] \dots ]]\bigr\|,
\\
&L^k_\Lambda(b) := \|\delta^k_\Lambda(b)\| 
                = \bigl\|[D_\Lambda, [D_\Lambda, \dots [D_\Lambda, b] \dots ]]\bigr\|.
\end{align*}

where the commutator is taken $k$ times.  
If these Lipschitz seminorms are actually Lipschitz norms , then we can consider the quantum 
Gromov--Hausdorff distance between the corresponding compact quantum 
metric spaces.


\subsection{Compact Quantum Groups}

Quantum groups emerged in the 1980s as a framework for studying symmetries 
of noncommutative spaces and exactly solvable models in statistical mechanics 
and quantum field theory. Among them, compact quantum groups, introduced by 
Woronowicz~\cite{12} in the late 1980s, provide a natural generalization of 
compact topological groups to the noncommutative setting. They play a central 
role in noncommutative geometry, representation theory, and the theory of 
operator algebras. In this subsection we recall the basic definition and some 
well‑known facts that will be needed later.

\begin{definition}\label{def:compact-quantum-group}
A \emph{compact quantum group} is a pair $(\mathcal{A}, \Delta)$ where 
$\mathcal{A}$ is a unital $C^*$-algebra and 
$\Delta : \mathcal{A} \to \mathcal{A} \otimes \mathcal{A}$ is a unital 
$*$-homomorphism, called the \emph{comultiplication}, satisfying
\begin{enumerate}
    \item $(\Delta \otimes \operatorname{id}_{\mathcal{A}}) \Delta 
          = (\operatorname{id}_{\mathcal{A}} \otimes \Delta) \Delta$ 
          as homomorphisms $\mathcal{A} \to \mathcal{A} \otimes \mathcal{A} 
          \otimes \mathcal{A}$ (coassociativity);
    \item the spaces 
          $(\mathcal{A} \otimes 1_{\mathcal{A}}) \Delta(\mathcal{A}) 
          := \operatorname{span}\{(a \otimes 1_{\mathcal{A}}) \Delta(b) : 
          a, b \in \mathcal{A}\}$ and 
          $(1_{\mathcal{A}} \otimes \mathcal{A}) \Delta(\mathcal{A})$ 
          are dense in $\mathcal{A} \otimes \mathcal{A}$ 
          (the cancellation property).
\end{enumerate}
Here the tensor product is the minimal $C^*$-tensor product.
\end{definition}

A fundamental fact about compact quantum groups is the existence and 
uniqueness of a \emph{Haar state} $h : \mathcal{A} \to \mathbb{C}$. 
It is a state satisfying the left and right invariance conditions
\[
(\operatorname{id}_{\mathcal{A}} \otimes h) \circ \Delta 
= h(\cdot) 1_{\mathcal{A}} 
= (h \otimes \operatorname{id}_{\mathcal{A}}) \circ \Delta.
\]
In general, the Haar state may not be a trace, and it may not be faithful 
on $\mathcal{A}$.

Two important classes of compact quantum groups are distinguished by 
additional properties of the Haar state. A compact quantum group 
$(\mathcal{A}, \Delta)$ is said to be of \emph{Kac type} if its Haar state 
is a trace, i.e.\ $h(ab) = h(ba)$ for all $a,b \in \mathcal{A}$. 
Equivalently, there exists a bounded antipode $S$ on a dense 
$*$-subalgebra satisfying $S^2 = \operatorname{id}$.

A compact quantum group is called \emph{coamenable} if its dual discrete 
quantum group is amenable. For a discrete group $\Gamma$, the compact 
quantum group $C^*(\Gamma)$ is coamenable precisely when $\Gamma$ is an 
amenable group. In the general theory, coamenability can be characterised 
by faithfulness of the Haar state on the reduced $C^*$-algebra, or by the 
existence of a bounded approximate invariant mean for the dual action.

\begin{example}[Discrete group duals]
Let \(\Gamma\) be a discrete group. Denote by \(C^*(\Gamma)\) the full group C$^*$-algebra. There is a natural comultiplication \(\Delta : C^*(\Gamma) \to C^*(\Gamma) \otimes C^*(\Gamma)\) determined by \(\Delta(\delta_g) = \delta_g \otimes \delta_g\) for the canonical unitaries \(\delta_g\) (\(g \in \Gamma\)). Together with the usual involution and unit, this makes \((C^*(\Gamma), \Delta)\) a compact quantum group. Its Haar state is given by the trivial character \(h(\lambda_g) = \delta_{g,e}\) for all \(g\), i.e. the \(*\)-homomorphism \(\varepsilon : C^*(\Gamma) \to \mathbb{C}\). It can be regarded as the dual of the discrete group \(\Gamma\).
\end{example}

Let $(\mathcal{H},\Lambda)$ denote the GNS representation of $C(\mathbb{G})$ 
associated with the Haar state, where $\mathcal{H}=L^{2}(\mathbb{G})$ and 
$\Lambda:C(\mathbb{G})\to L^{2}(\mathbb{G})$ is the canonical GNS map.

\subsubsection{The fundamental unitaries}
\label{subsubsec:fundamental_unitaries}

On the Hilbert space $L^{2}(\mathbb{G})\hat{\otimes}L^{2}(\mathbb{G})$, 
the \emph{left fundamental unitary} $W$ and the \emph{right fundamental 
unitary} $V$ are defined on the dense subspace 
$\Lambda\otimes\Lambda\bigl(C(\mathbb{G})\odot C(\mathbb{G})\bigr)$ by
\begin{align*}
    W^{*}\bigl(\Lambda(x)\otimes\Lambda(y)\bigr) 
      &= \Lambda\otimes\Lambda\bigl(\Delta(y)(x\otimes 1)\bigr), \\
    V\bigl(\Lambda(x)\otimes\Lambda(y)\bigr) 
      &= \Lambda\otimes\Lambda\bigl(\Delta(x)(1\otimes y)\bigr),
\end{align*}
for all $x,y\in C(\mathbb{G})$. Both unitaries implement the comultiplication 
of $C(\mathbb{G})$ via the fundamental relations
\begin{equation}
    \Delta(x) = W^{*}(1\otimes x)W = V(x\otimes 1)V^{*} 
    \qquad (x\in C(\mathbb{G})). \tag{2.3}
\end{equation}
Consequently, $\Delta$ extends uniquely to a normal $*$-homomorphism
\[
\Delta\colon L^{\infty}(\mathbb{G}) \longrightarrow 
L^{\infty}(\mathbb{G})\,\hat{\otimes}\,L^{\infty}(\mathbb{G}),
\]
turning $L^{\infty}(\mathbb{G})$ into a compact von Neumann algebraic quantum 
group~\cite{12,26}.

\subsubsection{Flip operators and leg-numbering notation}
\label{subsubsec:flips_notation}

The Hilbert space $L^{2}(\mathbb{G})\hat{\otimes}L^{2}(\mathbb{G})$ carries 
a natural \emph{flip unitary} $\Sigma$, defined on elementary tensors by
\[
\Sigma\bigl(\Lambda(x)\otimes\Lambda(y)\bigr) = \Lambda(y)\otimes\Lambda(x), 
\qquad x,y\in C(\mathbb{G}).
\]
Conjugation by $\Sigma$ implements the \emph{flip automorphism} $\sigma$ 
at the $C^{*}$-algebraic level:
\[
\sigma(a\otimes b) \;:=\; b\otimes a 
\;=\; \Sigma\,(a\otimes b)\,\Sigma, 
\qquad a,b\in\mathbb{B}\bigl(L^{2}(\mathbb{G})\bigr).
\]

We shall consistently use the following conventions:
\begin{enumerate}
    \item \emph{Leg‑numbering.} For an operator 
    $T\in\mathbb{B}\bigl(L^{2}(\mathbb{G})\hat{\otimes}L^{2}(\mathbb{G})\bigr)$ 
    we set
    \[
    T_{12} = T\otimes 1,\qquad
    T_{13} = (1\otimes\sigma)(T\otimes 1),\qquad
    T_{23} = 1\otimes T,
    \]
    and analogously for more legs. The same notation is employed for operators 
    on tensor products of more than two copies of $L^{2}(\mathbb{G})$.
    
    \item \emph{Sweedler notation.} Whenever an element $a\in C(\mathbb{G})$ 
    satisfies $\Delta(a)\in C(\mathbb{G})\odot C(\mathbb{G})$, we write
    \[
    \Delta(a) = a_{(0)}\otimes a_{(1)}.
    \]
    The same convention is used for more general coactions of $\mathbb{G}$; 
    sums are understood in the usual Sweedler fashion.
\end{enumerate}

Since $W$ and $\Sigma V \Sigma$ are unitary corepresentations and $W_{23}$ 
and $(\Sigma V\Sigma)_{13}$ commute, 
$Z = W \Sigma V \Sigma \in L^{\infty}(\mathbb{G})\bar{\otimes} 
\mathcal{B}\bigl(L^{2}(\mathbb{G})\bigr)$ is also a unitary corepresentation.

\subsubsection{Corepresentation theory}
\label{subsubsec:corep_theory}

A \emph{unitary corepresentation} of a compact quantum group $\mathbb{G}$ 
on a Hilbert space $H$ is a unitary element 
$U\in L^{\infty}(\mathbb{G})\,\bar{\otimes}\,\mathcal{B}(H)$ satisfying
\begin{equation}
    (\Delta\otimes\mathrm{id})(U)=U_{13}U_{23}. \tag{2.4}
\end{equation}
Both fundamental unitaries $W$ and $V$ satisfy the pentagon identity 
(for $W$ it reads $W_{12}W_{13}W_{23}=W_{23}W_{12}$); together with (2.4) 
this implies that $W$ and $\Sigma V\Sigma$ are unitary corepresentations 
of $\mathbb{G}$.

The corepresentation theory of $\mathbb{G}$ closely resembles the 
representation theory of a classical compact group. One has natural 
notions of direct sums, tensor products, equivalence and irreducibility. 
Moreover, every finite‑dimensional unitary corepresentation decomposes as 
a direct sum of irreducible ones. The unit element $1\in C(\mathbb{G})$ 
gives a one‑dimensional (irreducible) corepresentation, whose class is 
denoted by $e$.

For each $\alpha\in\Irred(\mathbb{G})$ the element $u^{\alpha}$ belongs to $L^{\infty}(\mathbb{G})\,\bar{\otimes}\,\mathcal{B}(H_{\alpha})$ for a 
finite dimensional Hilbert space $H_{\alpha}$; write $d_{\alpha}:=\dim H_{\alpha}$.  
Choosing an orthonormal basis of $H_{\alpha}$ we identify 
$L^{\infty}(\mathbb{G})\,\bar{\otimes}\,\mathcal{B}(H_{\alpha})
\cong M_{d_{\alpha}}\bigl(L^{\infty}(\mathbb{G})\bigr)$ and obtain the 
matrix coefficients $u_{ij}^{\alpha}\in L^{\infty}(\mathbb{G})$ 
($1\leq i,j\leq d_{\alpha}$). Actually one has $u_{ij}^{\alpha}\in C(\mathbb{G})$, 
and the linear span
\[
\Pol(\mathbb{G}):=\operatorname{span}_{\mathbb{C}}
\bigl\{\,u_{ij}^{\alpha}\mid \alpha\in\Irred(\mathbb{G}),\;1\leq i,j\leq d_{\alpha}\,\bigr\}
\]
forms a dense Hopf $*$‑subalgebra of $C(\mathbb{G})$, called the 
\emph{regular Hopf algebra} of $\mathbb{G}$. The comultiplication restricts 
to a map $\Delta:\Pol(\mathbb{G})\to\Pol(\mathbb{G})\odot\Pol(\mathbb{G})$ 
which on matrix coefficients reads
\begin{equation}
    \Delta(u_{ij}^{\alpha})=\sum_{k=1}^{d_{\alpha}} u_{ik}^{\alpha}\otimes u_{kj}^{\alpha}. \tag{2.5}
\end{equation}
The algebra $\Pol(\mathbb{G})$ also carries an antipode $S$ and a counit 
$\varepsilon$ (see, e.g., \cite[Theorem~5.4.1]{26}). Because 
$\Pol(\mathbb{G})$ is a Hopf $*$‑algebra, the operators $W^{*}$ and $V$ map 
$\Pol(\mathbb{G})\odot\Pol(\mathbb{G})$ bijectively onto itself 
\cite[Theorem~1.3.18]{26}; the same is true for $W$ and $V^{*}$.



\subsection{Fusion algebras}
Now we turn to the fusion algebra, a structure closely tied to the 
representation theory of quantum groups. 
Fusion algebras, also known as based rings, provide an algebraic framework 
that unifies the representation theory of compact groups and the combinatorial 
structure of discrete groups. These algebras first emerged systematically in 
the study of conformal field theory and have since found profound applications 
in subfactor theory, quantum groups, and tensor categories.

\begin{definition}\label{def:fusion_algebra}
An algebra $R = \bigoplus_{\xi \in I} \mathbb{C}\xi$ over $\mathbb{C}$ with a 
basis $I$ and an identity $e \in I$ is called a \emph{fusion algebra} if it 
satisfies the following axioms:

\begin{enumerate}
    \item[(i)] \textbf{Structure constants.} There exist non-negative integers 
    $N_{\xi,\eta}^{\zeta}$ for $\xi, \eta, \zeta \in I$ such that
    \[
    \xi\eta = \sum_{\zeta \in I} N_{\xi,\eta}^{\zeta}\,\zeta.
    \]
    These integers $N_{\xi,\eta}^{\zeta}$ are called the 
    \emph{structural constants} of $R$.
    
    \item[(ii)] \textbf{Conjugation.} There exists an involutive map 
    $\xi \mapsto \bar{\xi}$ on $I$, called \emph{conjugation}, whose 
    $\mathbb{Z}$-linear extension to $R$ satisfies
    \[
    \overline{\xi\eta} = \bar{\eta}\bar{\xi}, \quad \xi, \eta \in I.
    \]
    
    \item[(iii)] \textbf{Frobenius reciprocity.} The structural constants obey 
    the symmetry relations
    \[
    N_{\xi,\eta}^{\zeta} = N_{\xi,\zeta}^{\bar{\eta}} 
    = N_{\xi,\bar{\eta}}^{\bar{\zeta}}, \quad \xi, \eta, \zeta \in I.
    \]
    
    \item[(iv)] \textbf{Dimension function.} There exists a function 
    $d: I \to [1,\infty)$, called the \emph{dimension function} (or 
    \emph{Frobenius--Perron dimension}), such that 
    $d(\xi) = d(\bar{\xi})$ for all $\xi \in I$ and
    \[
    d(\xi)d(\eta) = \sum_{\zeta \in I} N_{\xi,\eta}^{\zeta}\,d(\zeta), 
    \quad \xi, \eta \in I.
    \]
\end{enumerate}
\end{definition}

\begin{remark}
The dimension function $d$ is uniquely determined by the above conditions 
and satisfies $d(e) = 1$. Moreover, the Frobenius reciprocity (iii) ensures 
that the algebra $R$ is semisimple when equipped with the bilinear form 
$\langle \xi, \eta \rangle = \delta_{\xi, \bar{\eta}}$.
\end{remark}

We now present two classical examples of fusion algebras: group algebras of 
discrete groups and representation rings of compact quantum groups.

\begin{example}\label{ex:group-algebra}
Let $G$ be a discrete group and $R = \mathbb{Z}[G] = \bigoplus_{g \in G} \mathbb{Z}g$ forms a fusion algebra 
with the following data:
\begin{itemize}
    \item Basis: $\Xi_0 = G$ with identity $e$ being the group identity;
    \item Multiplication: $N_{g_1,g_2}^{h} = \delta_{h,g_1g_2}$ for 
    $g_1, g_2, h \in G$;
    \item Conjugation: $\bar{g} = g^{-1}$;
    \item Dimension: $d(g) = 1$ for all $g \in G$.
\end{itemize}
This example represents the simplest class of fusion algebras, where all 
basis elements have trivial dimension and the structural constants are given 
by the group multiplication.
\end{example}

\begin{example}[Representation ring]\label{ex:rep-ring}
Let $\mathbb{G}$ be a compact matrix quantum group and let 
$\widehat{\mathbb{G}}$ denote the set of equivalence classes of irreducible 
unitary corepresentations of $\mathbb{G}$. The representation ring 
$\Xi = R(\mathbb{G}) = \bigoplus_{\xi \in \widehat{\mathbb{G}}} \mathbb{C}\xi$ 
carries a fusion algebra structure characterized by:
\begin{itemize}
    \item Basis: $I = \Irred(\mathbb{G})$, where $\Irred(\mathbb{G})$ is a 
    complete family of representatives for the equivalence classes of 
    irreducible unitary corepresentations of $\mathbb{G}$;
    \item Multiplication: $N_{\xi,\eta}^{\zeta}$ equals the multiplicity of 
    $\zeta$ in the irreducible decomposition of the tensor product 
    $\xi \otimes \eta$;
    \item Conjugation: $\bar{\xi}$ is the conjugate corepresentation of $\xi$;
    \item Dimension: $d(\xi) = \dim \xi$, the quantum dimension.
\end{itemize}
The dimension condition (iv) in Definition~\ref{def:fusion_algebra} reflects 
the well-known identity $\dim(\xi \otimes \eta) = (\dim \xi)(\dim \eta)$ in 
representation theory.
\end{example}

The notion of amenability for finitely generated fusion algebras was 
introduced in \cite{14}. A finitely generated fusion algebra is called 
\emph{amenable} if there exist $1 < p < \infty$ and a finitely supported, 
symmetric, non‑degenerate probability measure $\mu$ on $I$ such that
\[
\|\lambda_{p,\mu}\| = 1.
\]

If $\mathbb{G}$ is a compact matrix quantum group, its corepresentation ring 
$R(\mathbb{G})$ is finitely generated. Hence amenability of $\mathbb{G}$ can 
be defined naturally through the amenability of $R(\mathbb{G})$.

In \cite{3}, David Kyed proposed a more general definition that applies to 
quantum groups that do not necessarily possess a fundamental corepresentation. 
We shall introduce this definition in detail below.

\begin{definition}[Amenable Fusion Algebra]\label{def:amenable}
A fusion algebra $R = \mathbb{C}[I]$ is called \emph{amenable} if 
$1 \in \sigma(\lambda_{2,\mu})$ for every finitely supported, symmetric 
probability measure $\mu$ on $I$.
\end{definition}

Here $\sigma(\lambda_{2,\mu})$ denotes the spectrum of the operator 
$\lambda_{2,\mu}$. From Proposition~1.3 and Corollary~4.4 in \cite{14} 
it follows that the definition agrees with the one in \cite{14} on the 
class of finitely generated fusion algebras.

\begin{definition}[Boundary]\label{def:boundary}
Let $R = \mathbb{C}[I]$ be a fusion algebra. For two finite subsets 
$S, F \subseteq I$ we define the \emph{boundary of $F$ relative to $S$} as 
the set
\begin{align*}
\partial_S(F) &= \{\alpha \in F \mid \exists\, \xi \in S : 
\operatorname{supp}(\alpha\xi) \nsubseteq F\} \\
&\quad \cup \{\alpha \in F^c \mid \exists\, \xi \in S : 
\operatorname{supp}(\alpha\xi) \nsubseteq F^c\}.
\end{align*}
Here, and in what follows, $F^c$ denotes the set $I \setminus F$.
\end{definition}

\begin{definition}[F{\o}lner condition]\label{def:folner}
Let $R = \mathbb{C}[I]$ be a fusion algebra with dimension function $d$. 
For every finite, non-empty subset $S \subseteq I$ and every 
$\varepsilon > 0$ there exists a finite subset $F \subseteq I$ such that
\[
\sum_{\xi \in \partial_S(F)} d(\xi)^2 < \varepsilon 
\sum_{\xi \in F} d(\xi)^2.
\]
\end{definition}

\begin{remark}
By \cite{3}, a fusion algebra $R$ is amenable provided that it satisfies the 
F{\o}lner condition.
\end{remark}









\section{Compact Quantum Metric Space from Fusion Algebras}



Consider a compact quantum group $\mathbb{G}$ and leting $F(\mathbb{G})$ denoted the $\mathbb{C}$-vector space with basis $\operatorname{Irred}(\mathbb{G})$,we obtain a *-algebra called fusion algebra.
We equip $F(\mathbb{G})$ with the inner product for which
$\operatorname{Irred}(\mathbb{G})$ is an orthogonal basis, and denote by
$\ell^{2}(\operatorname{Irred}(\mathbb{G}))$ the completion.

Let $\{\alpha : \alpha \in \operatorname{Irred}(\mathbb{G})\}$ be the
canonical orthogonal basis of $\ell^{2}(\operatorname{Irred}(\mathbb{G}))$,
with
\[
\langle \alpha,\beta\rangle = \delta_{\alpha,\beta}\, d(\alpha)^2.
\]
Here $d:\operatorname{Irred}(\mathbb{G})\to [1,\infty)$ is the dimension
function satisfying
\[
d(\alpha)=d(\bar{\alpha}), \qquad
d(\alpha)d(\beta)=\sum_{\gamma\in \operatorname{Irred}(\mathbb{G})}
N_{\alpha,\beta}^{\gamma}\, d(\gamma),
\qquad \alpha,\beta\in \operatorname{Irred}(\mathbb{G}).
\]
We denote by $N_{\alpha,\beta}^{\gamma}$ the multiplicity of
$\gamma \in \operatorname{Irred}(\mathbb{G})$ in the decomposition of
$\alpha \otimes \beta$ into irreducibles; that is,
\begin{equation}
\alpha \otimes \beta \simeq
\bigoplus_{\gamma \in \operatorname{Irred}(\mathbb{G})}
\gamma^{\oplus N_{\alpha,\beta}^{\gamma}}.
\tag{3.1}
\end{equation}
Let $c_{c}(\operatorname{Irred}(\mathbb{G}))$ denote the dense subspace of
$\ell^{2}(\operatorname{Irred}(\mathbb{G}))$ spanned by this basis.

For each $f \in F(\mathbb{G})$, left multiplication defines a linear map
\[
\Lambda_{0}(f) \colon c_{c}(\operatorname{Irred}(\mathbb{G}))
\longrightarrow c_{c}(\operatorname{Irred}(\mathbb{G})),
\]
and by \cite[Lemma~3.3]{16}, $\Lambda_{0}(f)$ extends to a bounded operator
$\Lambda(f)$ on $\ell^{2}(\operatorname{Irred}(\mathbb{G}))$.


\begin{definition}\label{def:length}
Let $\mathbb{G}$ be a compact quantum group. A \emph{length function} on 
$\mathbb{G}$ is a function 
$\ell \colon \operatorname{Irred}(\mathbb{G}) \to [0,\infty)$ such that
\begin{enumerate}
    \item[(i)] $\ell(e) = 0$ ;
    \item[(ii)] $\ell(\bar{\alpha}) = \ell(\alpha)$ for all 
    $\alpha \in \operatorname{Irred}(\mathbb{G})$;
    \item[(iii)] $\ell(\gamma) \leq \ell(\alpha) + \ell(\beta)$ for all 
    $\alpha, \beta, \gamma \in \operatorname{Irred}(\mathbb{G})$ such that 
    $\gamma$ is equivalent to a sub-corepresentation of 
    $\alpha \otimes \beta$.
\end{enumerate}
\end{definition}

A length function $\ell$ is called \emph{proper} if 
$\ell^{-1}([0,N])$ is finite for all $N \geq 0$ and 
$\ell(\alpha) = 0$ only when $\alpha = e$.
The length function also defines an unbounded operator 
\[
\tilde{D}_\ell : c_c(\operatorname{Irred}(\mathbb{G})) \longrightarrow 
\ell^2(\operatorname{Irred}(\mathbb{G})),\qquad 
\tilde{D}_\ell(\alpha) = \ell(\alpha)\,\alpha.
\]
Assume moreover that $\ell$ is integer-valued.
We now additionally assume that $\mathbb{G}$ is equipped with a proper 
length function $\ell \colon \operatorname{Irred}(\mathbb{G}) \to 
\mathbb{N}$. Setting 
\[
A_{n} := \operatorname{span}\{\alpha \in \operatorname{Irred}(\mathbb{G}) 
\mid \ell(\alpha) \leqslant n\},
\]
then $A_m \subset A_n$ if $m < n$, 
$F(\mathbb{G}) = \bigcup_{n=0}^{\infty} A_n$, $A_n^* = A_n$, 
$A_m A_n \subseteq A_{m+n}$, and $A_0 = \mathbb{C}e$ where $e$ is the unit element. Consequently, 
the fusion algebra $F(\mathbb{G})$ belongs to the class of filtered 
$*$-algebras studied by Ozawa and Rieffel in \cite{14}.

We can view each $A_n$ as a finite-dimensional, hence closed, subspace 
of $\mathcal{H}$. Let $P_n$ denote the orthogonal projection of 
$\mathcal{H}$ onto $A_n$. We set $Q_0 = P_0$ and $Q_n = P_n - P_{n-1}$ for $n \ge 1$. Then, for each $a \in F(\mathbb{G})$, we define
\[
a_n = Q_n(a),
\]
so that $\tilde{D}_\ell(a_n) = n a_n$. In particular, if there is no $\alpha \in \operatorname{Irred}(\mathbb{G})$ such that $\ell(\alpha)=n$, then $a_n = 0$.
 The $Q_n$ are mutually orthogonal, and 
$\sum_{n=0}^\infty Q_n = I_{\mathcal{H}}$ in the strong operator 
topology. Then $a_n \in A_n$, but $a_n \notin A_{n-1}$ unless 
$a_n = 0$. Moreover, $a = \sum_n a_n$, with at most $p$ non‑zero terms 
if $a \in A_p$.

Actually, $\tilde{D}_\ell = \sum_{n=1}^\infty n Q_n$. We now consider 
whether it can form a spectral triple for the fusion algebra.

\begin{lemma}
Let $\mathbb{G}$ be a compact quantum group equipped with a proper 
length function $\ell$. Then the triple 
$\bigl(F(\mathbb{G}), \ell^2(\operatorname{Irred}(\mathbb{G})), 
\tilde{D}_\ell\bigr)$ is a spectral triple.
\end{lemma}
\begin{proof}
For all $x, y \in \ell^2(\operatorname{Irred}(\mathbb{G}))$,
\[
\langle \tilde{D}_\ell x, y \rangle 
= \sum_{\alpha} \overline{\ell(\alpha)}\,\overline{x_\alpha}\,y_\alpha
= \sum_{\alpha} \overline{x_\alpha}\,\ell(\alpha)\,y_\alpha 
= \langle x, \tilde{D}_\ell y \rangle,
\]
so $\tilde{D}_\ell$ is self‑adjoint.

We show that $\tilde{D}_\ell$ has compact resolvent. If 
$\operatorname{Irred}(\mathbb{G})$ is finite this is trivial. Assume 
$\operatorname{Irred}(\mathbb{G})$ is infinite and take 
$\lambda \notin \sigma(\tilde{D}_\ell)$. The resolvent 
$(\tilde{D}_\ell - \lambda)^{-1}$ is the diagonal operator with 
eigenvalues $(\ell(\alpha) - \lambda)^{-1}$. Since $\ell$ is proper, 
the set
\[
\ell^{-1}([0,R]) = \{\alpha \in \operatorname{Irred}(\mathbb{G}) : 
\ell(\alpha) \le R\}
\]
is finite for every $R \ge 0$. Consequently, the eigenvalues of 
$(\tilde{D}_\ell - \lambda)^{-1}$ tend to zero, i.e.\ the resolvent 
is compact.

Now suppose $a \in A_p$. For any $m, n \ge 0$, if 
$Q_m a Q_n \neq 0$ then there exists $\xi \in A_n$ with 
$a\xi \in A_m$. Because $A_p A_n \subseteq A_{p+n}$, we have 
$p+n \ge m$. Taking adjoints, $Q_m a^* Q_n \neq 0$ implies 
$p+m \ge n$. Hence $|m-n| \le p$. Therefore
\[
a = \sum_{|m-n| \le p} Q_m a Q_n,
\]
converging in the strong operator topology. For each $j$ with 
$|j| \le p$ set
\[
T_j = \sum_m Q_m a Q_{m-j}.
\]
For fixed $j$, the ranges of the terms $Q_m a Q_{m-j}$ are mutually 
orthogonal, as are their domains, so
\[
\|T_j\| = \sup_m \|Q_m a Q_{m-j}\| \le \|a\|.
\]
For any $m, n$ we have 
$[\tilde{D}_\ell, Q_m a Q_n] = (m-n)Q_m a Q_n$. In particular,
\[
[\tilde{D}_\ell, Q_m a Q_{m-j}] = j\, Q_m a Q_{m-j},
\]
and thus $[\tilde{D}_\ell, T_j] = j T_j$. Since 
$a = \sum_{|j|\le p} T_j$, we obtain
\[
[\tilde{D}_\ell, a] = \sum_{|j|\le p} j T_j,
\]
which is a bounded operator. Therefore 
$\bigl(F(\mathbb{G}), \ell^2(\operatorname{Irred}(\mathbb{G})), 
\tilde{D}_\ell\bigr)$ is a spectral triple.
\end{proof}
When considering a discrete group equipped with a length function, one often studies the properties of rapid decay and polynomial growth. In \cite{20}, Vergnioux introduced a notion of rapid decay for discrete quantum groups. We will apply this property of rapid decay to fusion algebras constructed from the irreducible corepresentations of a compact quantum group. Since every $f \in F(\mathbb{G})$ has finite support, we can directly give a definition on the fusion algebra.


\begin{definition}[Rapid Decay for Fusion Algebras]\label{rapid decay}
Let $\mathbb{G}$ be a compact  quantum group.Fix a length function \(\ell : \operatorname{Irred}(\mathbb{G}) \to [0,\infty)\). For \(f \in F(\mathbb{G})\) define the weighted \(\ell^2\)-norm
\[
\|f\|_{2,s} := \left( \sum_{\alpha \in \operatorname{Irred}(\mathbb{G})} |f(\alpha)|^2 \, d(\alpha)^2 \, (1+\ell(\alpha))^{2s} \right)^{1/2} \qquad s>0.
\]
 The fusion algebra \(F(\mathbb{G})\) is said to have the \emph{rapid decay property} if there exist constants \(C>0\) and \(s>0\) such that
\[
\|\Lambda(f)\| \le C \|f\|_{2,s}
\]
for every \(f \in F(\mathbb{G})\).
\end{definition}

\begin{definition}\label{polynomial growth}
    Let \( G \) be a compact quantum group. We say that \( G \) has 
    \emph{polynomial growth of order s with respect to length function \(\ell\) on \( G \)} 
    if there exist \(c\in\mathbb{R}_+\) such that
    \[
    \sum_{\substack{\ell(\alpha) \in (n-1, n]}} d(\alpha)^2 \leq c n^s
    \]
    for all \( n \in \mathbb{N} \).
    We say that \( G \) has \emph{strong polynomial growth of order s with respect to 
    length function \(\ell\) on \( G \)} if there exist 
    \(c_1,c_2\in\mathbb{R}_+\) such that
    \[
    c_2 n^s \leq \sum_{\substack{\ell(\alpha) \in (n-1, n]}} d(\alpha)^2 
    \leq c_1 n^s
    \]
    for all \( n \in \mathbb{N} \).

    We say that \( G \) has \emph{(strong) polynomial growth} if it has 
    (strong) polynomial growth with respect to some length function on \( G \).
\end{definition}


\begin{proposition}\label{prposition rapid decay}
    Let \(G\) be a compact  quantum group with polynomial growth. 
    Then \(\mathcal{F}(\mathbb{G})\) is rapid decay.
\end{proposition}
\begin{proof}
    For any unit vector 
    \(\xi = \sum_{\beta} \xi_\beta \beta\) in 
    \(\ell^2(\operatorname{Irred}(\mathbb{G}))\), 
    \(f=\sum_\alpha f(\alpha)\alpha\in\mathcal{F}(\mathbb{G})\) we have
    \begin{align*}
        \|\Lambda(f)\|
        &= \sup_{\xi}\Bigl\| \sum_\alpha \sum_\beta f(\alpha)\xi(\beta) 
           \sum_\gamma N_{\alpha,\beta}^{\gamma} \gamma \Bigr\| \\
        &\leq \sum_\alpha |f(\alpha)| 
               \sup_{\xi}\Bigl\| \sum_\gamma \sum_\beta N_{\alpha,\beta}^{\gamma} 
               \xi_\beta \gamma \Bigr\| \\
        &= \sum_\alpha |f(\alpha)| 
               \sup_\xi\Bigl( \sum_\gamma \Bigl( \sum_\beta N_{\alpha,\beta}^{\gamma} 
               \xi_\beta \Bigr)^2 d(\gamma)^2 \Bigr)^{\frac{1}{2}} \\
        &\leq \sum_\alpha |f(\alpha)| 
                \sup_\xi\Bigl( \sum_\gamma \sum_\beta (N_{\alpha,\beta}^{\gamma} 
               \xi_\beta)^2 d(\alpha)^2 d(\gamma)^2  
               \Bigr)^{\frac{1}{2}} \\
        &\leq \sum_\alpha |f(\alpha)| 
                \sup_\xi\Bigl( \sum_\beta (\xi_\beta)^2 \Bigl( \sum_\gamma 
               N_{\alpha,\beta}^{\gamma} d(\gamma) \Bigr)^2 d(\alpha)^2 
               \Bigr)^{\frac{1}{2}} \\
        &= \sum_\alpha |f(\alpha)| 
            \sup_\xi\Bigl( \sum_\beta (\xi_\beta)^2 d(\beta)^2 d(\alpha)^4 
           \Bigr)^{\frac{1}{2}} \\
        &= \sum_\alpha |f(\alpha)| d(\alpha)^2 \\
        &\leq \Bigl( \sum_{\alpha} \bigl( |f(\alpha)|^2 d(\alpha)^2 
           (1+\ell(\alpha))^{2s} \bigr) \Bigr)^{\frac{1}{2}}
           \Bigl( \sum_{\alpha} \frac{d(\alpha)^2}{(1+\ell(\alpha))^{2s}} 
           \Bigr)^{\frac{1}{2}} .
    \end{align*}
    Since \(G\) has polynomial growth,
    \(C = \bigl( \sum_{\alpha} \frac{d(\alpha)^2}{(1+\ell(\alpha))^{2s}} 
    \bigr)^{\frac{1}{2}}\) is a constant. We complete the proof.
\end{proof}



Now we begin to consider Lip-seminorm on fusion algebra.
If we define \(\delta(f):=[\tilde{D}_\ell,\Lambda(f)]\), then 
\(\delta^k(f)= [\tilde{D}_\ell, [\tilde{D}_\ell, \dots [\tilde{D}_\ell,\Lambda(f)]\dots ]]\),
where the commutator is taken \(k\) times.

\begin{lemma}
    For the fusion algebra \(F(\mathbb{G})\) and for every positive integer 
    \(k\), \(L_k(f)=\| \delta^k(f)\|
\) is a Lip-seminorm.
\end{lemma}
\begin{proof}
    For any \(f = \sum_{\alpha} f(\alpha)\,\alpha \in F(\mathbb{G})\), let 
    \(\Lambda(f)\) denote its left regular representation on 
    \(\ell^2(\operatorname{Irred}(\mathbb{G}))\):
    \[
    \Lambda(f)\beta = \sum_{\alpha,\gamma\in\operatorname{Irred}(\mathbb{G})} 
    f(\alpha) N_{\alpha,\beta}^{\gamma} \,\gamma.
    \]

    The adjoint \([\tilde{D}_\ell,\Lambda(f)]^*\) satisfies
    \[
    \langle \beta, [\tilde{D}_\ell,\Lambda(f)]^* \gamma \rangle
    = \overline{\langle \gamma, [\tilde{D}_\ell,\Lambda(f)]\beta \rangle}
    = \sum_{\alpha} (\ell(\gamma)-\ell(\beta)) \overline{f(\alpha)} 
      N_{\alpha,\beta}^{\gamma}.
    \]

    And 
    \begin{align*}
    \langle \beta, [\tilde{D}_\ell,\Lambda(f^*)] \gamma \rangle
    &= (\ell(\beta)-\ell(\gamma)) \langle \beta,\Lambda(f^*)\gamma \rangle \\
    &= (\ell(\beta)-\ell(\gamma)) \sum_{\alpha} f^*(\alpha) 
       N_{\bar\alpha,\gamma}^{\beta} \\
    &= \sum_{\alpha} (\ell(\beta)-\ell(\gamma)) \overline{f(\alpha)} 
       N_{\bar\alpha,\gamma}^{\beta}.
    \end{align*}

    For a compact quantum group, the fusion coefficients satisfy 
    \(N_{\bar\alpha,\gamma}^{\beta}=N_{\alpha,\beta}^{\gamma}\) 
    (the representation category is spherical). Using this,
    \[
    \langle \beta, [\tilde{D}_\ell,\Lambda(f^*)] \gamma \rangle
    = \sum_{\alpha} (\ell(\beta)-\ell(\gamma)) \overline{f(\alpha)} 
      N_{\alpha,\beta}^{\gamma}
    = -\sum_{\alpha} (\ell(\gamma)-\ell(\beta)) \overline{f(\alpha)} 
      N_{\alpha,\beta}^{\gamma}.
    \]

    Comparing with the matrix elements of \([\tilde{D}_\ell,\Lambda(f)]^*\) we obtain
    \[
    [\tilde{D}_\ell,\Lambda(f)]^* = -\,[\tilde{D}_\ell,\Lambda(f^*)].
    \]
    If we define \(\delta(f):=[\tilde{D}_\ell,\Lambda(f)]\), then 
    \(\delta^k(f)= [\tilde{D}_\ell, [\tilde{D}_\ell, \dots 
    [\tilde{D}_\ell,\Lambda(f)]\dots ]]\).
    Now, for \(k \in \mathbb{N}\) and \(f \in  F(\mathbb{G})\), consider
    \[
    \delta^k(f)^* = (-1)^k \delta^k(f^*) .
    \]
    Then \(\|\delta^k(f)^*\| = \|\delta^k(f^*)\|\).
    Note that if \(e\in\operatorname{Irred}(\mathbb{G})\) is the identity of 
    \(F(\mathbb{G})\), 
    \(\delta^k(f)e = \sum_{\alpha\in\operatorname{Irred}(\mathbb{G})} 
    \ell(\alpha)^k f(\alpha)\alpha\).
    Then for \(f\in F(\mathbb{G})\), we have \(\|\delta^k(f)\|=0\) if and only if 
    \(f\in \mathbb{C}e\).
    So the Lipschitz semi‑norm (Rieffel's framework) can be defined as the 
    operator norm of this commutator:
    \[
    L_k(\cdot) := \|\delta^k(\cdot)\|.
    \]
\end{proof}

\begin{remark}
    In the truncated picture, for \(\Lambda>0\) let \(P_\Lambda\) be the 
    orthogonal projection onto 
    \(A_\Lambda = \operatorname{span}\{\alpha: \ell(\alpha)\le\Lambda\}\),
    set \(\tilde{D}_{\ell\Lambda}=P_\Lambda\tilde{D}_\ell P_\Lambda\), and for 
    \(f_\Lambda\in P_\Lambda F(\mathbb{G})P_\Lambda\) (which corresponds to a 
    function \(f\) on \(\operatorname{Irred}(\mathbb{G})\) via 
    \(f_\Lambda = P_\Lambda\Lambda(f)P_\Lambda\)) we have
    \[
    [\tilde{D}_{\ell\Lambda},f_\Lambda]\beta =
    \sum_\gamma \sum_{\substack{\alpha:\ \exists\beta,\gamma\in A_\Lambda \\
    \text{s.t. } N_{\alpha,\beta}^{\gamma}\neq 0}}
    (\ell(\gamma)-\ell(\beta)) f(\alpha) N_{\alpha,\beta}^{\gamma} \,\gamma,
    \qquad \gamma,\beta\in A_\Lambda.
    \]
    Note that for \(f\in F(\mathbb{G})\), we have 
    \([\tilde{D}_{\ell\Lambda},f_\Lambda]=0\) iff \(f_\Lambda\in \mathbb{C}e\).
    and 
    \[
    [\tilde{D}_{\ell\Lambda},\Lambda(f_\Lambda)]^* = -\,[\tilde{D}_{\ell\Lambda},\Lambda(f_\Lambda^*)].
    \]
    
   Let $\delta_\Lambda^k(\cdot) =
    [ \tilde{D}_{\ell\Lambda},
      [ \tilde{D}_{\ell\Lambda},
        \dots [ \tilde{D}_{\ell\Lambda},
          \Lambda(\cdot)]
         \dots ] ]$,then
 $\|\delta_\Lambda^k(\cdot)\|$ is a \(\operatorname{Lip}\)-seminorm for \(A_\Lambda\). Since \(A_\Lambda\) 
    is finite dimensional, \(\|\delta_\Lambda^k(\cdot)\|\) is a Lip-norm.
\end{remark}

For any Lip seminorm \(\|\cdot\|\) on \(F(\mathbb{G})\), \(f \in F(\mathbb{G})\) 
and \(r \geqslant 0\) we denote the corresponding closed balls as follows:
\[
\mathbb{B}_r^{\|\cdot\|}(f) := \{g \in F(\mathbb{G}) \mid \|f - g\| \leqslant r\} .
\]
Here we let \([\cdot]: F(\mathbb{G}) \to F(\mathbb{G})/\mathbb{C}\) denote the 
quotient map and \(\|\cdot\|_{F(\mathbb{G})/\mathbb{C}}\) denote the quotient 
norm on \(F(\mathbb{G})/\mathbb{C}\).

\begin{theorem}\label{compact quantum metric space on fusion algebra}
    If \(\mathbb{G}\) is a compact quantum group with polynomial growth of order s, then 
    for any positive integer \(k>s\), the algebra \(F(\mathbb{G})\) together 
    with the seminorm \(L_\ell^k\) is a compact quantum metric space.
\end{theorem}
\begin{proof}
    By Rieffel~\cite{9}, we know \((F(\mathbb{G}), L_\ell^k)\) is a compact quantum 
    metric space if and only if the subset 
    \(\overline{\mathbb{B}}_1^{L_\ell^k}(0)\subseteq F(\mathbb{G})/\mathbb{C}\) is 
    totally bounded with respect to the quotient norm 
    \(\|\cdot\|_{F(\mathbb{G})/\mathbb{C}}\) on \(F(\mathbb{G})/\mathbb{C}\). 
    Let \(B=\overline{\mathbb{B}}_1^{L_\ell^k}(0)=B_1+B_2\) where
    \[
    B_1 = \Bigl\{ f \in B \Bigm| f = \sum_{\alpha\,|\,\ell(\alpha)\le n} 
          f(\alpha)\alpha \Bigr\}
    \]
    and
    \[
    B_2 = \Bigl\{ f \in B \Bigm| f = \sum_{\alpha\,|\,\ell(\alpha)> n} 
          f(\alpha)\alpha \Bigr\}.
    \]
    For all \(f \in B_1\),
    \begin{align*}
        \|f\| &\le \sum_{\ell(\alpha)\le n} |f(\alpha)| d(\alpha)^2 \\
        &\le \biggl( \sum_{\ell(\alpha)\le n} |f(\alpha)|^2 d(\alpha)^2 
           (1+\ell(\alpha))^{2s}) \biggr)^{1/2}
           \biggl( \sum_{\ell(\alpha)\le n} 
           \frac{d(\alpha)^2}{(1+\ell(\alpha))^{2s}} \biggr)^{1/2} \\
        &\le C_n .
    \end{align*}
    where 
    \( C_n = 2^s\Bigl( \sum_{\ell(\alpha)\le n} 
      \frac{d(\alpha)^2}{(1+\ell(\alpha))^{2s}} \Bigr)^{1/2} \).\\
    For all \(f \in B_2\), fix a positive integer \(k> s\). By the rapid 
    decay property we have
    \begin{align*}
        \|f\| &\le C \Bigl( \sum_{\ell(\alpha)> n} |f(\alpha)|^2 d(\alpha)^2 
               (1+\ell(\alpha)^{2s}) \Bigr)^{1/2} \\
        &\le C 2^s n^{s-k} \Bigl( \sum_{\ell(\alpha)> n} |f(\alpha)|^2 
           d(\alpha)^2 (\ell(\alpha))^{2k} \Bigr)^{1/2}.
    \end{align*}
    Then \(B_1\) is a bounded subset of a finite‑dimensional normed space, and 
    thus totally bounded. Moreover, by our choice of \(n\), the set \(B_2\) is 
    contained in the \(\varepsilon\)-ball around \(0\). Therefore, for any positive 
    integer \(k> s\), the algebra \(F(\mathbb{G})\) together with the 
    seminorm \(L_\ell^k\) is a compact quantum metric space.
\end{proof}

\begin{example}[Connected compact Lie groups]
    Let \(G\) be a connected compact Lie group and equip the commutative
    $C^*$ algebra \(C(G)\) with a length function
    \(\ell(\alpha)=\|\lambda_\alpha\|\), where \(\lambda_\alpha\) is the highest 
    weight of an irreducible representation \(\alpha\).
    Then \(d(\alpha)\) is the ordinary dimension of the representation.
    The Weyl dimension formula yields
    \[
    \sum_{\ell(\alpha)\in(n-1,n]} d(\alpha)^2 \asymp n^{D-1},
    \]
    where \(D\) is the dimension of \(G\) as a smooth manifold.
    Thus \(C(G)\) has strong polynomial growth of order \(D\).
    For instance, for all \(N\in\mathbb{N}\), \(\mathrm{SU}(N)\) and 
    \(\mathrm{SO}(N)\) have strong polynomial growth.
\end{example}












\section{Gromov--Hausdorff Convergence of Spectral
Truncation for Compact Quantum Groups}

Let $\mathbb{G}$ be a compact quantum group with Hopf algebra of polynomial 
functions $\operatorname{Pol}(\mathbb{G})$. Let $\widehat{\mathbb{G}}$ be the 
set of equivalence classes of irreducible unitary representations.
Let $\{u_{ij}^\alpha\}$ be the matrix coefficients of 
$\alpha \in \widehat{\mathbb{G}}$, and let $(\mathcal{H}, \Lambda)$ denote the 
GNS representation of $C(\mathbb{G})$ associated with the Haar state. Then 
$\{\Lambda(u_{ij}^\alpha)\}$ forms an orthonormal basis of $\mathcal{H}$.
According to Definition~\ref{def:length}, for a proper length function, an 
unbounded Dirac operator 
$D_\ell:\Lambda(\operatorname{Pol}(\mathbb{G}))\to \mathcal{H}$ can be defined by
\[
D_\ell\bigl(\Lambda(u_{i,j}^\alpha)\bigr) = \ell(\alpha)\,\Lambda(u_{i,j}^\alpha) .
\]
By \cite[Lemma~5.4]{23}, we know that 
$(\operatorname{Pol}(\mathbb{G}),\mathcal{H},D_\ell)$ is a spectral triple. 

A spectral truncation for the spectral triple is given by a spectral 
projection $P_\Lambda$ of $D_\ell$ onto the eigenspaces with eigenvalues of 
modulus $\leq\Lambda$. The $C^*$-algebra $C(\mathbb{G})$ is replaced by 
$P_\Lambda C(\mathbb{G}) P_\Lambda$, acting on the Hilbert space 
$P_\Lambda \mathcal{H}$, and we set 
$D_{\ell,\Lambda}:=D_\ell|_{P_\Lambda\mathcal{H}}$. 
If $L_{C(\mathbb{G})}$ is the Lipschitz seminorm on $C(\mathbb{G})$ associated with $D_\ell$, we 
consider the quantum Gromov--Hausdorff distance between $C(\mathbb{G})$ and 
its truncation. If this distance tends to zero as $\Lambda\to\infty$, we say 
that the spectral truncation of $C(\mathbb{G})$ converges in quantum 
Gromov--Hausdorff distance. 

When we consider the spectral truncation of this Dirac operator, since the 
length function is proper the resulting spectral truncation is analogous to 
the classical Peter--Weyl truncation.
The commutator seminorm 
$L_\ell^k: \operatorname{Pol}(\mathbb{G}) \to [0,\infty)$ is defined by
\[
L_\ell^k(a) := \|[\cdots[D_\ell,[D_\ell, a]]\cdots]\|,
\]
where $[D_\ell, a] = D_\ell a - a D_\ell$ , the commutator is taken $k$ times and $\|\cdot\|$ is the operator norm 
on $L^2(\mathbb{G})$.

We now prove that $L_\ell^k$ is a Lip-seminorm in the sense of Rieffel.


\begin{lemma}\label{lemma:lip-norm}
    The seminorm $L_\ell^k$ is a Lip-seminorm on $C(\mathbb{G})$.
\end{lemma}

\begin{proof}
    Let
    \[
        \delta^k(a):=\underbrace{[D_\ell,[D_\ell,\dots,[D_\ell,a]\dots]]}_{k\text{ times}}
    \]
    denote the $k$-fold iterated commutator with $D_\ell$.

    Since $(\operatorname{Pol}(\mathbb{G}),L^2(\mathbb{G}),D_\ell)$ is a spectral triple,
    the seminorm $L_\ell$ is densely defined on $\operatorname{Pol}(\mathbb{G})$.
    For $k=1$, we have
    \[
        \delta(u_{ij}^\alpha)\Lambda(1)
        =[D_\ell,\Lambda(u_{ij}^\alpha)]\Lambda(1)
        =D_\ell \Lambda(u_{ij}^\alpha)-\Lambda(u_{ij}^\alpha)D_\ell \Lambda(1)
        =D_\ell \Lambda(u_{ij}^\alpha).
    \]
    Assume now that
    \[
        \delta^k(u_{ij}^\alpha)\Lambda(1)=D_\ell^k \Lambda(u_{ij}^\alpha).
    \]
    Then
    \begin{align*}
        \delta^{k+1}(u_{ij}^\alpha)\Lambda(1)
        &= [D_\ell,\delta^k(u_{ij}^\alpha)]\Lambda(1) \\
        &= D_\ell\,\delta^k(u_{ij}^\alpha)\Lambda(1)-\delta^k(u_{ij}^\alpha)D_\ell\Lambda(1) \\
        &= D_\ell^{k+1}\Lambda(u_{ij}^\alpha).
    \end{align*}
    Hence, by induction,
    \[
        \delta^k(u_{ij}^\alpha)\Lambda(1)=D_\ell^k \Lambda(u_{ij}^\alpha)
        \qquad \text{for all } k\geq 1.
    \]

    Now let
    \[
        a=\sum_{\alpha,i,j} a_{ij}^\alpha\, \Lambda(u_{ij}^\alpha)\in C(\mathbb{G}).
    \]
    If $L_\ell^k(a)=0$, then $\|\delta^k(a)\|=0$. Therefore,
    \[
        0=\delta^k(a)\Lambda(1)
        =\sum_{\alpha,i,j} a_{ij}^\alpha\, \delta^k(u_{ij}^\alpha)\Lambda(1)
        =\sum_{\alpha,i,j} a_{ij}^\alpha\, D_\ell^k\Lambda(u_{ij}^\alpha).
    \]
    Since the matrix coefficients $\Lambda(u_{ij}^\alpha)$ are linearly independent, it follows that
    \[
        \ell(\alpha)^k a_{ij}^\alpha=0
        \qquad \text{for all } \alpha,i,j.
    \]
    By the definition of the length function $\ell$, we have $\ell(\alpha)=0$ only for the trivial representation $\alpha=1$.
    Hence $a_{ij}^\alpha=0$ whenever $\alpha\neq 1$, so $a$ must be a scalar multiple of the identity.

    Therefore, the kernel of $L_\ell^k$ consists exactly of the scalar multiples of the identity.
    Since $L_\ell^k(a^*)=L_\ell^k(a)$ is immediate, we conclude that $L_\ell^k$ is a Lip-seminorm on $C(\mathbb{G})$.
\end{proof}

\begin{remark}
    For any $k \in \mathbb{N}$ and $a \in \operatorname{Pol}(\mathbb{G})$, 
    define the iterated commutator
    \[
    \delta^k(a) = [D_\ell, [D_\ell, \dots [D_\ell, a]\dots]] \quad 
    (k\text{ times}),
    \]
    and set
    \[
    L_\ell^k(a) = \|\delta^k(a)\|.
    \]
    Analogously, for $b \in P_\Lambda \operatorname{Pol}(\mathbb{G}) P_\Lambda$ let
    \[
    \delta_\Lambda^k(b) = [D_{\ell,\Lambda}, [D_{\ell,\Lambda}, \dots 
                         [D_{\ell,\Lambda}, b]\dots]]
    \]
    and
    \[
    L_{\ell,\Lambda}^k(b) = \|\delta_\Lambda^k(b)\|.
    \]
    The same reasoning as in the proof of Lemma~\ref{lemma:lip-norm} shows 
    that $L_\ell^k$ is a Lip‑seminorm on $\mathbb{G}$, while its truncation 
    $L_{\ell,\Lambda}^k$ is a Lip‑norm on the finite‑dimensional subspace 
    $P_\Lambda C(\mathbb{G}) P_\Lambda$.
\end{remark}

\begin{definition}
    A Lip-seminorm $L$ on a compact quantum group $(\mathcal{A}, \Delta)$ is 
    said to be \emph{left-invariant} (or \emph{right-invariant}) if for all 
    $a \in \mathcal{A}$ with $a^* = a$ and all $\mu \in S(\mathcal{A})$ we have
    \[
    L(a * \mu) \leq L(a) \quad (\text{or } L(\mu * a) \leq L(a)),
    \]
    where $a * \mu = (\mu \otimes \mathrm{id}) \Delta(a)$ and similarly for 
    $\mu * a$. Then $L$ is said to be \emph{bi-invariant} if it is both left 
    and right invariant.
\end{definition}

Now we show that for any $k\in \mathbb{N}$, $L_\ell^k$ is left-invariant. 
First, the comultiplication extends to a normal $*$-homomorphism 
$\Delta \colon L^{\infty}(\mathbb{G}) \to 
L^{\infty}(\mathbb{G}) \bar{\otimes} L^{\infty}(\mathbb{G})$, turning 
$L^{\infty}(\mathbb{G})$ into a compact von Neumann algebraic quantum 
group~\cite{22}. On $L^2(\mathbb{G}) \hat{\otimes} L^2(\mathbb{G})$ we have the 
left fundamental unitary $W$ and right fundamental unitary $V$ defined on the dense subspace 
$\Lambda \otimes \Lambda (C(\mathbb{G}) \odot C(\mathbb{G}))$ by the relation
\[
W^*(\Lambda(x) \otimes \Lambda(y)) = \Lambda \otimes \Lambda\bigl(\Delta(y)(x \otimes 1)\bigr),
\]
\[V(\Lambda(x) \otimes \Lambda(y)) = \Lambda \otimes \Lambda\bigl(\Delta(x)(1\otimes y)\bigr).\]
It implements the comultiplication on $C(\mathbb{G})$ by means of
\[
\Delta(x) = W^*(1 \otimes x)W=V(x\otimes1)V^*. \tag{2.3}
\]

\begin{lemma}\label{lemma left}
\(L_\ell^k\) is bi-invariant, i.e.\ the seminorm \(L_\ell^k\) satisfies
\[
L_\ell^k\bigl((\varphi \otimes \operatorname{id})\Delta(a)\bigr) 
\le \|\varphi\| \cdot L_\ell^k(a),
\]
\[L_\ell^k\bigl(( \operatorname{id}\otimes\varphi)\Delta(a)\bigr) 
\le \|\varphi\| \cdot L_\ell^k(a).
\]
for all \(a \in \operatorname{Pol}(\mathbb{G})\) and 
\(\varphi \in S(C(\mathbb{G}))\).
\end{lemma}

\begin{proof}
Computing with matrix coefficients we obtain
\begin{align*}
(\mathrm{id} \otimes D_{\ell}) W^{*}\bigl(\Lambda(u_{ij}^{\alpha}) \otimes \Lambda(u_{pq}^{\beta})\bigr) 
&= (\mathrm{id} \otimes D_{\ell})(\Lambda \otimes \Lambda)
   \bigl(\Delta(u_{pq}^{\beta})(u_{ij}^{\alpha} \otimes 1)\bigr) \\
&= (\mathrm{id} \otimes D_{\ell})\Bigl( \sum_{k=1}^{d_{\beta}} 
   \Lambda(u_{pk}^{\beta} u_{ij}^{\alpha}) \otimes \Lambda(u_{kq}^{\beta}) \Bigr) \\
&= \ell(\beta) \sum_{k=1}^{d_{\beta}} \Lambda(u_{pk}^{\beta} u_{ij}^{\alpha}) \otimes \Lambda(u_{kq}^{\beta}) \\
&= W^{*}(\mathrm{id} \otimes D_{\ell}) \bigl(\Lambda(u_{ij}^{\alpha}) \otimes \Lambda(u_{pq}^{\beta})\bigr).
\end{align*}

For all \(a \in \operatorname{Pol}(\mathbb{G})\), we compute as follows, where all 
equalities are understood to hold on the dense subspace 
\(\Lambda(\operatorname{Pol}(\mathbb{G})) \odot \Lambda(\operatorname{Pol}(\mathbb{G}))
\subset L^2(\mathbb{G}) \hat\otimes L^2(\mathbb{G})\):
\[
\begin{aligned}
&[\operatorname{id} \otimes D_\ell,\,
 [\operatorname{id} \otimes D_\ell,\,
  \dots,\,
 [\operatorname{id} \otimes D_\ell, \Delta(a)]\dots]] \\
&\quad = [\operatorname{id} \otimes D_\ell,\,
 [\operatorname{id} \otimes D_\ell,\,
  \dots,\,
 [\operatorname{id} \otimes D_\ell, W^*(1\otimes a)W]\dots]] \\
&\quad = W^*\,[\operatorname{id} \otimes D_\ell,\,
 [\operatorname{id} \otimes D_\ell,\,
  \dots,\,
 [\operatorname{id} \otimes D_\ell, 1\otimes a]\dots]]\,W.
\end{aligned}
\]

 
For all \(\varphi \in S(C(\mathbb{G}))\),we  obtain the following identities for operators defined on the 
dense subspace \(\operatorname{Pol}(\mathbb{G}) \subset L^2(\mathbb{G})\):
\begin{align*}
L_\ell^k\bigl((\varphi \otimes \operatorname{id})\Delta(a)\bigr)
&= \bigl\| [D_\ell, [D_\ell, \dots, [D_\ell, 
   (\varphi \otimes \operatorname{id})\Delta(a)]\dots]] \bigr\| \\
&= \bigl\| \varphi(a_{(0)}) \, [D_\ell, [D_\ell, \dots, [D_\ell, a_{(1)}]\dots]] \bigr\| \\
&= \bigl\| (\varphi \otimes \operatorname{id}) 
   [\operatorname{id} \otimes D_\ell,\,
    [\operatorname{id} \otimes D_\ell,\,
     \dots,\,
    [\operatorname{id} \otimes D_\ell, \Delta(a)]\dots]] \bigr\| \\
&= \bigl\| (\varphi \otimes \operatorname{id}) 
   \bigl( W^* (1 \otimes [D_\ell, [D_\ell, \dots, [D_\ell, a]\dots]]) W \bigr) \bigr\| \\
&\le \|\varphi \otimes \operatorname{id}\| \cdot 
   \bigl\| W^* (1 \otimes [D_\ell, [D_\ell, \dots, [D_\ell, a]\dots]]) W \bigr\| \\
&= \|\varphi\| \cdot L_\ell^k(a).
\end{align*}
The same argument applies to the right-invariance case, and we obtain  \[L_\ell^k\bigl(( \operatorname{id}\otimes\varphi)\Delta(a)\bigr) 
\le \|\varphi\| \cdot L_\ell^k(a).\]

Thus the desired inequality holds.
\end{proof}

\begin{remark}
As in Lemma~\ref{lemma left}, the same reasoning shows that the seminorm 
\(L_{\ell,\Lambda}^k\) remains bi-invariant on the truncated space. Indeed, 
\(L_{\ell,\Lambda}^k\) is defined on the compressed spectral subspace 
\(P_\Lambda \mathcal{H}\) via the same iterated commutator formula involving the 
truncated Dirac operator \(D_{\ell,\Lambda}\). Consequently, the proof of 
Lemma~\ref{lemma left} carries over verbatim to the truncated setting. Hence 
\(L_{\ell,\Lambda}^k\) is bi-invariant on \(C(\mathbb{G})_\Lambda\): for all \(b \in P_\Lambda \operatorname{Pol}(\mathbb{G}) P_\Lambda\) 
and \(\varphi \in S(P_\Lambda C(\mathbb{G}) P_\Lambda)\), we have
\[
L_{\ell,\Lambda}^k\bigl((\varphi \otimes \operatorname{id})\Delta(b)\bigr) 
\le \|\varphi\| \cdot L_{\ell,\Lambda}^k(b),
\]
\[
L_{\ell,\Lambda}^k\bigl(( \operatorname{id}\otimes\varphi)\Delta(b)\bigr) 
\le \|\varphi\| \cdot L_{\ell,\Lambda}^k(b).
\]
\end{remark}

\begin{lemma}\label{lem:approx_state}
Let \(\bigl(C(\mathbb{G}), L_\ell^k\bigr)\) be a compact quantum metric space. 
Then for any state \(\mu \in S(C(\mathbb{G}))\) and any \(\varepsilon > 0\), there 
exists a state \(\nu \in S(C(\mathbb{G}))\) satisfying
\[
|\mu(a) - \nu(a)| \;\leq\; \varepsilon \, L_\ell^k(a), \qquad \forall a \in C(\mathbb{G}),
\]
such that \(\nu(b) = 0\) whenever 
\(b \in \bigl(P_\Lambda C(\mathbb{G})\bigr)^c\).
\end{lemma}

\begin{proof}
Let \((\mathcal{H}, \Lambda)\) denote the GNS representation of \(C(\mathbb{G})\) 
associated with the Haar state. Since \((C(\mathbb{G}), \Delta)\) is in reduced 
form, the representation \(\Lambda\) is faithful. Define
\[
\mathcal{K} \;:=\; \Lambda\bigl(\operatorname{Pol}(\mathbb{G})\bigr)\,\mathcal{H},
\]
where \(\operatorname{Pol}(\mathbb{G})\) is the dense \(*\)-subalgebra of 
\(C(\mathbb{G})\) consisting of matrix coefficients. Then \(\mathcal{K}\) is 
a dense subspace of \(\mathcal{H}\).

Fix a state \(\mu \in S(C(\mathbb{G}))\) and \(\varepsilon > 0\). By 
\cite[Proposition~4.7]{11}, there exists a finite convex combination 
\(\nu\) of vector states, where the vectors are taken from \(\mathcal{K}\), 
such that
\[
|\mu(a) - \nu(a)| \;\leq\; \varepsilon \, L_\ell^k(a), \qquad \forall a \in C(\mathbb{G}).
\]

We now construct the required support condition for \(\nu\). Let 
\(a_1, \dots, a_m \in \operatorname{Pol}(\mathbb{G})\) be  elements 
in \(\mathcal{K}\) that determine the vector states comprising \(\nu\). For 
each \(j = 1, \dots, m\), define \(F_{a_j}\) as in \cite[Proposition~2.2]{11}. 
Set
\[
F \;:=\; \bigcup_{j=1}^m F_{a_j}.
\]

By \cite[Proposition~2.2]{11}, the set \(F\) consists of finitely many 
irreducible representations of \(\mathbb{G}\), and the vector states 
corresponding to \(a_j\) are supported on the finite-dimensional subspaces 
\(P_F \mathcal{H}\), where \(P_F\) denotes the projection onto the isotypic 
component of \(F\). Consequently, the state \(\nu\) is supported on 
\(P_F \mathcal{H}\), which implies
\[
\nu(b) = 0 \quad \text{for all } b \in \bigl(P_F C(\mathbb{G}) \bigr)^c.
\]
Since \(F\) is finite, we have \(P_F \leq P_\Lambda\) for some \(\Lambda\). 
Consequently, \(P_F C(\mathbb{G})  \subseteq P_\Lambda C(\mathbb{G})\), 
which yields the desired property.
\end{proof}


\begin{theorem}\label{theorem GH}
Let \(\mathbb{G}\) be a coamenable compact quantum group equipped with a proper length 
function. If \(\bigl(C(\mathbb{G}), L_\ell^k\bigr)\) is a compact quantum 
metric space, then there exists \(\Lambda\) such that we have quantum 
Gromov--Hausdorff convergence
\[
\operatorname{dist}_q\bigl((C(\mathbb{G}),L^k_{\ell}),\, 
(P_\Lambda C(\mathbb{G}), L^k_{\ell,\Lambda})\bigr) \to 0.
\]
\end{theorem}

\begin{proof}
Let \(\epsilon\) be the coidentity of \(\mathbb{G}\), viewed as an element 
of \(S(C(\mathbb{G}))\). By Lemma~\ref{lem:approx_state}, there exists a 
sequence \(\chi_\Lambda\) converging to \(I^{\mathbb{G}}\) in the weak\(^*\) 
topology on \(S(C(\mathbb{G}))\).

Define an operator \(P_{\chi_\Lambda}\) on \(P_\Lambda C(\mathbb{G})\) by
\[
P_{\chi_\Lambda}(u) = \chi_\Lambda * u = (\operatorname{id} \otimes \chi_\Lambda)\Delta(u), 
\qquad u\in C(\mathbb{G}).
\]
The range of \(P_{\chi_\Lambda}\) is contained in \(C(\mathbb{G})_\Lambda\), 
and \(L_{\ell,\Lambda}^k(P_{\chi_\Lambda}(u)) \le L_\ell^k(u)\) follows directly 
from the left‑invariance property and the definition of \(P_{\chi_\Lambda}(u)\).

Now for every self‑adjoint \(u\in C(\mathbb{G})\) (i.e.\ \(u=u^*\)) we have
\[
\begin{aligned}
\|u - P_{\chi_\Lambda}(u)\|
   &= \|\epsilon * (u - P_{\chi_\Lambda}(u))\| \\
   &= \sup_{\phi \in S(\mathbb{G})} \bigl|\phi\bigl(\epsilon *(u - \chi_\Lambda * u)\bigr)\bigr| \\
   &= \sup_{\phi \in S(\mathbb{G})} \bigl|\epsilon(u * \phi) - \chi_\Lambda(u * \phi)\bigr| \\
   &\le \sup_{\phi \in S(\mathbb{G})} \varepsilon\, L_\ell^k(u * \phi) \\
   &\le \varepsilon\, L_\ell^k(u),
\end{aligned}
\]
where the last inequality uses the left‑invariance of \(L_\ell^k\).

Applying Proposition~8.5 of \cite{9} immediately yields that for every 
\(\varepsilon>0\),
\[
\operatorname{dist}_q\bigl((C(\mathbb{G}),L_{\ell}),\, 
(P_\Lambda C(\mathbb{G}), L_{\ell,\Lambda})\bigr) \to 0
\quad\text{for all } \Lambda \ge N.
\]
Thus the required convergence holds.
\end{proof}

However, when considering spectral truncation, the truncation space is given by 
$C(\mathbb{G})_\Lambda := P_\Lambda C(\mathbb{G}) P_\Lambda$. 
In this setting, Proposition 8.5 of \cite{9} no longer holds. 
Consequently, we must  develop a new approach to address this issue.
\begin{definition}[Right Coaction]
A \emph{right coaction} $\alpha$ of the function algebra $A$ on the operator system $X$ is a uci map $\alpha: X \to X \otimes A$ such that the \emph{coaction property}
\begin{equation}\label{eq:coaction-prop}
(\alpha \otimes \mathbf{I}^A)\alpha = (\mathbf{I}^X \otimes \Delta)\alpha
\end{equation}
and the Podle\'s density condition
\[
\operatorname{span}((\mathbf{1}_X \otimes A)\alpha(X)) = X \otimes A
\]
are satisfied. A \emph{left coaction} $\beta: X \to A \otimes X$ is defined analogously.
\end{definition}

For a coaction \(\alpha\colon X\to X\otimes A\), an element \(x\in X\) is called a 
\emph{fixed point} if it satisfies \(x_{(0)}\otimes x_{(1)} = x\otimes \mathbf{1}_A\), 
and the set of all fixed points is denoted by 
\(X^\alpha := \{x\in X \mid \alpha(x)=x\otimes \mathbf{1}_A\}\). 
The coaction \(\alpha\) is said to be \emph{ergodic} if its only fixed points are 
scalar multiples of the unit, i.e.\ \(X^\alpha = \mathbb{C}\mathbf{1}_X\). 
For example, the comultiplication \(\Delta\colon A\to A\otimes A\) is ergodic.

\begin{definition}
We denote by $\tau_{\Lambda}: \mathcal{B}(H) \to \mathcal{B}(H_{\Lambda})$ the \emph{compression map}, given by
\[
\tau_{\Lambda}(T) := P_{\Lambda} T P_{\Lambda},
\]
for all $T \in \mathcal{B}(H)$, and write $A_{\Lambda} := \tau_{\Lambda}(A) \subseteq \mathcal{B}(H_{\Lambda})$ for the image of the function algebra $A$ under the compression map.
\end{definition}

\begin{remark}
Throughout this section we may drop the subindex $\Lambda$ of the projection $P_{\Lambda}$ and the compression map $\tau_{\Lambda}$ whenever convenient.
Note that $A_{\Lambda}$ is an operator system and the compression map $\tau: A \to A_{\Lambda}$ is ucp onto.By \cite{6},We know there exist the following conclusion: there exists unique ergodic commuting right and left coactions 
\(\alpha^{\tau}\colon A_{\Lambda}\to A_{\Lambda}\otimes A\) and \(\beta^{\tau}\colon A_{\Lambda}\to A\otimes A_{\Lambda}\) 
such that 
\((\tau\otimes \mathbf{I}^A)\Delta= \alpha^{\tau}\tau\) and \(( \mathbf{I}^A\otimes\tau)\Delta= \beta^{\tau}\tau\). 
Using Sweedler notation,the equation  reads
\[
\tau(x_{(0)})\otimes x_{(1)} = (\tau(x))_{(0)}\otimes (\tau(x))_{(1)} 
= y_{(0)}\otimes y_{(1)}\in A_{\Lambda}\otimes A,
\]
for all \(x\in A\) and \(y\in A_{\Lambda}\) with \(\tau(x)=y\).

Since the coaction \(\Delta\) is ergodic, then \(\alpha^{\tau}\) is ergodic as well.The same argument applies to the left coaction.
Assume that the function algebra $A$ is equipped with a seminorm $L_A$. For all elements $x \in A_\Lambda$, set
\begin{equation}
\begin{aligned}
L_{ A_\Lambda}^{\alpha^\tau}(x) &:= \sup_{\phi \in S(A_\Lambda)} L_A\!\bigl(\phi(x_{(0)})\,x_{(1)}\bigr), \\
L_{ A_\Lambda}^{\beta^\tau}(x) &:= \sup_{\phi \in S(A_\Lambda)} L_A\!\bigl(x_{(0)}\,\phi(x_{(1)})\bigr), \\
L_{A_{\Lambda}}^{\alpha^\tau,\beta^\tau}
&:= \max\!\left\{L_{A_{\Lambda}}^{\alpha^\tau},\, L_{A_{\Lambda}}^{\beta^\tau}\right\}.
\end{aligned}
\end{equation}
Since $\alpha^\tau$ is unital and $L_\ell^k$ is a regular Lipschitz seminorm, the induced seminorms
$L_{C(\mathbb{G})_\Lambda}^{\alpha^\tau}$, $L_{C(\mathbb{G})_\Lambda}^{\beta^\tau}$, and
$L_{C(\mathbb{G})_\Lambda}^{\alpha^\tau,\beta^\tau}$ are Lip-norms.

 \end{remark}


 


\begin{definition}
Let $\phi \in \mathcal{S}(A_{\Lambda})$ be any state. We denote the associated slice map by $\sigma_\Lambda^\phi : A_{\Lambda} \to A$, i.e.,
\[
\sigma_\Lambda^\phi(x) := \phi(x_{(0)}) x_{(1)} = (\phi \otimes \mathbf{I}^A) \alpha^\tau(x),
\]
for all $x \in A_{\Lambda}$. We call the map $\sigma_\Lambda^\phi$ a \emph{symbol map}.
\end{definition}

Firstly, we compute the compositions of the compression and symbol maps:
\begin{equation}
\sigma^\phi \tau(a) = \phi(\tau(a)_{(0)}) \tau(a)_{(1)} = \tau^*\phi(a_{(0)}) a_{(1)},
\end{equation}
for all $a \in A$, and
\begin{equation}
\tau \sigma^\phi(x) = \phi(x_{(0)}) \tau(x_{(1)}) = \tau^*\phi(a_{(0)}) \tau(a_{(1)}),
\end{equation}
for all $x \in A_{\Lambda}$ and $a \in A$ with $\tau(a) = x$ where $\tau^*$.


\begin{lemma}
    The maps $\tau: C(\mathbb{G}) \to C(\mathbb{G})_\Lambda$ and $\sigma^\phi : C(\mathbb{G})_\Lambda \to C(\mathbb{G})$ are unital and positive. Furthermore, for all $a \in C(\mathbb{G})$ and $x \in C(\mathbb{G})_\Lambda$, we have
    \[
        L_{C(\mathbb{G})_\Lambda}^{\alpha^{\tau},\beta^{\tau}}(\tau(a)) \leq L_\ell^k(a)
        \quad\text{and}\quad
        L_\ell^k(\sigma^\phi(x)) \leq L_{C(\mathbb{G})_\Lambda}^{\alpha^{\tau},\beta^{\tau}}(x).
    \]
\end{lemma}

\begin{proof}
    The map $\tau$ is clearly unital and positive. The map $\sigma^\phi$ is the composition of the unital completely positive maps $\alpha^\tau$ and $\phi \otimes \mathbf{I}^A$; hence it is also unital and positive.
    
    For  $a \in C(\mathbb{G})$ we compute:
    \begin{align*}
        L_{C(\mathbb{G})_\Lambda}^{\alpha^{\tau}}(\tau(a))
        &= \sup_{\phi \in \mathcal{S}(C(\mathbb{G})_\Lambda)} L_{\ell}^k\bigl(\phi(\tau(a)_{(0)}) \tau(a)_{(1)}\bigr) \\
        &= \sup_{\phi \in \mathcal{S}(C(\mathbb{G})_\Lambda)} L_{\ell}^k\bigl(\tau^* \phi(a_{(0)}) a_{(1)}\bigr) \\
        &= \sup_{\psi \in \tau^*\mathcal{S}(C(\mathbb{G})_\Lambda)} L_{\ell}^k\bigl(\psi(a_{(0)}) a_{(1)}\bigr) \\
        &\leqslant \sup_{\psi \in \mathcal{S}(C(\mathbb{G}))} L_{\ell}^k\bigl(\psi(a_{(0)}) a_{(1)}\bigr)\\
        &\leqslant L_{\ell}^k(a),
    \end{align*}
    where the inequality follows from the inclusion $\tau^*\mathcal{S}(C(\mathbb{G})_\Lambda) \subseteq \mathcal{S}(C(\mathbb{G}))$, and the last equality uses the left invariance of $L_\ell^k$. The inequality for $\beta^\tau$ is analogous. Consequently,
    \[
        L_{C(\mathbb{G})_\Lambda}^{\alpha^{\tau},\beta^{\tau}}(\tau(a)) \leq L_\ell^k(a).
    \]
    
    Conversely, for $x \in C(\mathbb{G})_\Lambda$, we have
    \begin{align*}
        L_\ell^k(\sigma^\phi(x))
        &= L_\ell^k(\phi(x_{(0)}) x_{(1)}) \\
        &\leq \sup_{\phi \in \mathcal{S}(C(\mathbb{G})_\Lambda)} L_\ell^k(\phi(x_{(0)}) x_{(1)}) \\
        &= L_{C(\mathbb{G})_\Lambda}^{\alpha^{\tau}}(x).
    \end{align*}
    This completes the proof.
\end{proof}

By \cite{6} proposition 6.8,we know that if \( L_\ell^k \) is a Lipschitz norm on \( C(\mathbb{G}) \) with \( \ker(L_\ell^k) = \mathbb{C}1_{ C(\mathbb{G})} \), and $L_{C(\mathbb{G})_\Lambda}^\alpha$  can be  the induced seminorm on $C(\mathbb{G})$. Let \( \mu, \nu \in S(C(\mathbb{G})) \) be states on \( A \) and consider the induced slice maps \( C(\mathbb{G})_\Lambda \to C(\mathbb{G})_\Lambda \), given by \( x \mapsto x_{(0)}\mu(x_{(1)}) \) and \( x \mapsto x_{(0)}\nu(x_{(1)}) \) respectively. Then the following holds, for all \( x \in C(\mathbb{G})_\Lambda \):
\[
\| x_{(0)}\mu(x_{(1)}) - x_{(0)}\nu(x_{(1)}) \| \leq 2 \, d^{L_\ell^k}(\mu, \nu) \, L_{C(\mathbb{G})_\Lambda}^\alpha(x).
\]

Similarly, if \( \beta: {C(\mathbb{G})_\Lambda} \to C(\mathbb{G}) \otimes {C(\mathbb{G})_\Lambda} \) is a left coaction and \( L_{C(\mathbb{G})_\Lambda}^\beta \) the induced seminorm on \( {C(\mathbb{G})_\Lambda} \), the following holds, for all \( x \in {C(\mathbb{G})_\Lambda} \):
\[
\| \mu(x_{(-1)}) \, x_{(0)} - \nu(x_{(-1)}) \, x_{(0)} \| \leq 2 \, d^{L_\ell^k}(\mu, \nu) \, L_{C(\mathbb{G})_\Lambda}^\beta(x).
\]
 If we consider $a\in C(\mathbb{G})$,since $L_\ell^k$ is bi-invariant,we have the follow includion:
 \[
\| a_{(0)}\mu(a_{(1)}) - a_{(0)}\nu(a_{(1)}) \| \leq 2 \, d^{L_\ell^k}(\mu, \nu) \, L_\ell^k(a),
\]
and
\[
\| \mu(a_{(-1)}) \, a_{(0)} - \nu(a_{(-1)}) \, a_{(0)} \| \leq 2 \, d^{L_\ell^k}(\mu, \nu) \, L_\ell^k(a).
\]



\begin{theorem}\label{theorem Gromov Hausdorff}
Let \(\mathbb{G}\) be a compact quantum group equipped with a proper length 
function. If \(\bigl(C(\mathbb{G}), L_\ell^k\bigr)\) is a compact quantum 
metric space, then  we have quantum 
Gromov--Hausdorff convergence
\[
\operatorname{dist}_q\bigl((C(\mathbb{G}),L_{\ell}^k),\, 
(C(\mathbb{G})_\Lambda, L_{C(\mathbb{G})_\Lambda}^{\alpha^{\tau},\beta^{\tau}})\bigr) \to 0 \qquad as \quad\Lambda\to \infty.
\]
\end{theorem}
\begin{proof}
By the above discuss,we know that for all $a\in C(\mathbb{G}),x\in C(\mathbb{G})_\Lambda$ with $\tau(a)=x$,
$\|\sigma^{\phi}\tau(a)-a\|=\|\tau^*\phi(a_{(0)})a_{(1)}-\epsilon(a_{(0)})a_{(1)}\|\leq 2d^{L_\ell^k}(\tau^*\phi,\epsilon)L_\ell^k(a)$;and
\begin{align*}
    \|\tau\sigma^{\phi}(x)-x\|&=\|\tau^*\phi(a_{(0)})\tau(a_{(1)})-\epsilon(a_{(0)})\tau(a_{(1)})\|\\
    &\leq 2\sup_{\psi\in S(C(\mathbb{G}))_\Lambda}|\psi(\tau^*\phi(a_{(0)})\tau(a_{(1)})-\epsilon(a_{(0)})\tau(a_{(1)})|\\
    &=2\sup_{\psi\in S(C(\mathbb{G}))_\Lambda}|(\tau^*\phi-\epsilon)(a_{(0)}\tau^*\psi(a_{(1)})|\\
    &\leq 2\sup_{\psi\in S(C(\mathbb{G}))_\Lambda}d^{L_\ell^k}(\tau^*\phi,\epsilon)L_\ell^k(x_{(0)}\psi(x_{1}))\\
    &=2d^{L_\ell^k}(\tau^*\phi,\epsilon)L_{C(\mathbb{G})_\Lambda}^{\alpha^{\tau},\beta^{\tau}}.
\end{align*}
By \cite{6} corollary 7.12,we see that $\|\sigma^{\phi}\tau(a)-a\|\leq \varepsilon L_\ell^k(a)$ and $\|\tau\sigma^{\phi}(x)-x\|\leq \varepsilon L_{C(\mathbb{G})_\Lambda}^{\alpha^{\tau},\beta^{\tau}}$ for all $a\in C(\mathbb{G})$ and $x\in C(\mathbb{G})_\Lambda$.
Hence $\tau$ and $\sigma^\phi$ satisfy the two condition of \ref{prop:gromov hausdorff converges}.Then we have 
\[
\operatorname{dist}_q\bigl((C(\mathbb{G}),L_{\ell}^k),\, 
(C(\mathbb{G})_\Lambda, L_{C(\mathbb{G})_\Lambda}^{\alpha^{\tau},\beta^{\tau}})\bigr) \to 0 \qquad as \quad\Lambda\to \infty.
\]
\end{proof}





















\section{Gromov--Hausdorff Convergence of Spectral Truncation for  
Coamenable Compact Quantum Group of Kac Type and Discrete Quantum Groups}

In this section, we discuss two examples: coamenable compact quantum groups 
of Kac type and discrete quantum groups. 

For coamenable compact quantum groups of Kac type, it follows from \cite{16} 
that $(C(\mathbb{G}), L_\ell)$ is a compact quantum metric space if and only 
if $(F(\mathbb{G}), L_\ell)$ is a compact quantum metric space. Later in this 
section, we will show that $(C(\mathbb{G}), L_\ell^k)$ is a compact quantum 
metric space if and only if $(F(\mathbb{G}), L_\ell^k)$ is a compact quantum 
metric space. And, by the results of Theorem~\ref{compact quantum metric space}, 
if $\mathbb{G}$ has polynomial growth, then there exists $k$ such that $(C(\mathbb{G}), L_\ell^k)$ is a 
compact quantum metric space.

On the other hand, for discrete quantum groups, it is known from \cite{23} 
that a discrete quantum group of rapid decay can also be equipped with a 
compact quantum metric space structure.

We first recall the definition of the algebra of central functions,
\begin{equation}\label{eq:central-functions}
    C_z(\mathbb{G}) := \{a \in C(\mathbb{G}) \mid \sigma\Delta(a) = \Delta(a)\},
\end{equation}
where $\sigma$ denotes the flip automorphism on $C(\mathbb{G})\otimes C(\mathbb{G})$.
This algebra characterizes precisely the conditions under which 
$\bigl(C(\mathbb{G}),\, L_\ell \bigr)$, where $L_\ell$ denotes the Lip-seminorm 
induced by the length function $\ell$, forms a \emph{compact quantum metric space} 
in the sense of Rieffel.

Indeed, by \cite[Theorem~A]{16}, if $\mathbb{G}$ is a compact, coamenable 
quantum group of Kac type and $\ell:\operatorname{Irred}(\mathbb{G})\to[0,\infty)$ 
is a proper length function, then $(C(\mathbb{G}),L_\ell)$ is a compact quantum 
metric space if and only if $(C_z(\mathbb{G}),L_\ell)$ is.

Let $\Lambda$ be the GNS representation of $C(\mathbb{G})$ on $L^2(\mathbb{G})$. 
By construction, the unbounded operator $D_\ell$ on $L^2(\mathbb{G})$ (defined 
in Section~3.1) leaves the subspace $\Lambda(\operatorname{Pol}_z(\mathbb{G}))$ 
invariant. Its restriction to this subspace is essentially self‑adjoint; we 
denote its self‑adjoint closure (still written as $D_\ell^z$) by the same symbol. 
Consequently $(\operatorname{Pol}_z(\mathbb{G}), L^2_z(\mathbb{G}), D_\ell^z)$ 
becomes a spectral triple, where 
$L^2_z(\mathbb{G}) := \overline{\Lambda(\operatorname{Pol}_z(\mathbb{G}))}$.

\begin{definition}
    A state $\varphi\in S(C(\mathbb{G}))$ is called \emph{central} if 
    $(\varphi \otimes \mathrm{id})\Delta = (\mathrm{id}\otimes \varphi)\Delta$.
\end{definition}

\begin{remark}
    In \cite{16}, Austad and Kyed defined a normal conditional expectation 
    $E:\mathbb{B}(L^2(\mathbb{G}))\to \mathbb{B}(L^2(\mathbb{G}))$ given by 
    $E(a)=(h\otimes \mathrm{id})\partial(a)$, where 
    $\partial(\cdot):=Z^*(1\otimes \cdot)Z$. And they prove that a state 
    $\varphi\in S(C(\mathbb{G}))$ is central if and only if 
    $\varphi=\varphi \circ E$ and there exists a sequence of central states 
    $\{\chi_n\}_{n\in \mathbb{N}}$ where are supported in only finitely many matrix cofficients converging to $\epsilon$ in the weak$^*$ 
    topology on $S(C(\mathbb{G}))$.
\end{remark}

\begin{theorem}\label{compact quantum metric space}
 If $\mathbb{G}$ is a compact, coamenable quantum group of Kac type and 
 $\ell:\operatorname{Irred}(\mathbb{G})\to[0,\infty)$ is a proper length 
 function, then $(C(\mathbb{G}),L_\ell^k)$ is a compact quantum metric space 
 if and only if $(C_z(\mathbb{G}),L_\ell^k)$ is.
\end{theorem}

\begin{proof}
 Similarly to $W$, the operator $\mathrm{id} \otimes D_\ell$ commutes with 
 $\Sigma V \Sigma$, $\Sigma V^* \Sigma$, $Z$ and $Z^*$, as operators on the 
 dense subspace 
 $\Lambda(\operatorname{Pol}(\mathbb{G}))\odot\Lambda(\operatorname{Pol}(\mathbb{G}))
 \subset L^2(\mathbb{G})\hat{\otimes}L^2(\mathbb{G})$.

 For $a,x,y \in \operatorname{Pol}(\mathbb{G})$, consider the relation between 
 $L_\ell^k(E(a))$ and $L_\ell^k(a)$:
\begin{align*}
    &\langle \Lambda(x),\delta^k(E(a))\Lambda(y)\rangle\\
    &=\langle \Lambda(x),[D_\ell,[D_\ell,\dots,[D_\ell,E(a)]\dots]]\Lambda(y)\rangle \\
    &= \sum_{l=0}^{k}(-1)^{k-l}
       \langle\Lambda(x),D_\ell^l\bigl((h\otimes \mathrm{id})\partial(a)\bigr)
       D_\ell^{k-l}\Lambda(y)\rangle\\
    &=\sum_{l=0}^{k}(-1)^{k-l}
       \langle \Lambda(1)\otimes D_\ell^l\Lambda(x),\,Z^*(1\otimes a) Z(\Lambda(1)\otimes D_\ell^{k-l}\Lambda(y))\rangle\\
    &=\sum_{l=0}^{k}(-1)^{k-l}
       \langle \Lambda(1)\otimes\Lambda(x),Z^*
       (1\otimes D_\ell^l(a)D_\ell^{k-l})Z\,(\Lambda(1)\otimes \Lambda(y))\rangle\\
    &=\langle \Lambda(x),\,(h\otimes \mathrm{id})\partial(
       [D_\ell,[D_\ell,\dots,[D_\ell,a]\dots]])\,\Lambda(y)\rangle\\
    &=\langle \Lambda(x),E(\delta^k(a))\Lambda(y)\rangle .
\end{align*}

 Then $\delta^k(E(a)) = E(\delta^k(a))$ and by \cite[Proposition~4.2]{16} we 
 know the map defines a normal conditional expectation, hence 
 $L_\ell^k(E(a)) = \|\delta^k(E(a))\| = \|E(\delta^k(a))\| \le L_\ell^k(a)$.

 Since for any two central states $\varphi$ and $\psi$ on $C(\mathbb{G})$, 
 $\varphi = \varphi \circ E$ and $\psi = \psi \circ E$. When $\mathbb{G}$ is 
 of Kac type, by \cite[Proposition~4.2]{16} we know $E$ maps 
 $\operatorname{Pol}(\mathbb{G})$ to $\operatorname{Pol}_z(\mathbb{G})$. Then
\begin{align*}
    d_\ell(\varphi, \psi) &:= \sup \bigl\{ |\varphi(a) - \psi(a)| \; ; \; 
        a \in \operatorname{Pol}(\mathbb{G}), \; L_\ell^k(a) \leq 1 \bigr\} \\
    &= \sup \bigl\{ |\varphi(a) - \psi(a)| \; ; \; 
        a \in \operatorname{Pol}_z(\mathbb{G}), \; L_\ell^k(a) \leq 1 \bigr\} \\
    &=: d_\ell^z \bigl( \varphi|_{C_z(\mathbb{G})}, \; \psi|_{C_z(\mathbb{G})} \bigr).
\end{align*}

 For any $\chi\in S(C(\mathbb{G}))$ and for any $a\in C(\mathbb{G})$, we denote 
 by $\alpha_\chi:C(\mathbb{G})\to C(\mathbb{G})$ the ucp map defined by 
 $\alpha_\chi(a)=(\mathrm{id}\otimes \chi)\Delta(a)$.
\begin{align*}
    \|a-\alpha_\chi(a)\|
    &=\sup\{\,|\varphi(a-\alpha_\chi(a))| \mid \|\varphi\|\leq 1\,\}  \\
    &=\sup\{\,|\varphi(\mathrm{id}\otimes(\epsilon-\chi))\Delta a| 
        \mid \|\varphi\|\leq 1\,\}\\
    &\leq \sup\{\,d_\ell(\epsilon,\chi) \cdot L_\ell^k((\varphi\otimes\mathrm{id})\Delta a) 
        \mid \|\varphi\|\leq 1\,\}\\
    &\leq d_\ell(\epsilon,\chi)\cdot L_\ell^k(a).
\end{align*}

 By Remark~5.2 , there exists a sequence of central states 
 $\chi_n$ converging to $\epsilon$, hence we obtain:
\begin{align*}
    \|a-\alpha_{\chi_n}(a)\|
    &=\|\alpha_\epsilon(a)-\alpha_{\chi_n}(a)\| \\
    &\leq d_\ell(\epsilon,\chi_n)\cdot L_\ell^k(a)
     = d_\ell\bigl(\epsilon|_{C_z(\mathbb{G})},\chi_n|_{C_z(\mathbb{G})}\bigr)
       \cdot L_\ell^k(a).
\end{align*}

 Since each $\alpha_{\chi_n}$ is ucp with finite-dimensional image and 
 $\lim_{n\to\infty} d_\ell^z\bigl(\epsilon|_{C_z(\mathbb{G})},\chi_n|_{C_z(\mathbb{G})}\bigr)\\=0$, 
 by \cite[Theorem~3.1]{27} the proof is complete.
\end{proof}


In \cite{26}, Thomas Timmerman introduced the Schur orthogonality relations 
between the fusion algebra and the central function algebra. Furthermore, in 
\cite{16}, Austad and Kyed established several lemmas regarding these relations. 
Summarizing these results, we obtain the following proposition.

\begin{proposition}\label{prop unitary}
If $\mathbb{G}$ is a compact quantum group, then the spectral triples 
$\bigl(F(\mathbb{G}), \ell^2(\operatorname{Irred}(\mathbb{G})), \tilde{D}_\ell\bigr)$ 
and 
$\bigl(\operatorname{Pol}_z(\mathbb{G}), L^2_z(\mathbb{G}), D_\ell^z\bigr)$ 
are unitarily equivalent via the unitary $\tilde{\chi}$.
\end{proposition}

\begin{proof}
We construct the unitary $\tilde{\chi}$ explicitly and then verify that it 
intertwines the representations and the Dirac operators.

\emph{Step 1: Definition and unitary extension of $\tilde{\chi}$.}  
Define a map $\chi : c_c(\operatorname{Irred}(\mathbb{G})) \to L^2_z(\mathbb{G})$ 
on the basis elements by
\[
\chi(\alpha) = \sum_{i=1}^{d_\alpha} \Lambda(u_{ii}^\alpha),\qquad 
\alpha\in\operatorname{Irred}(\mathbb{G}),
\]
where $\Lambda:\operatorname{Pol}(\mathbb{G})\hookrightarrow L^2(\mathbb{G})$ is 
the canonical inclusion. By construction, the image of $\chi$ is dense in 
$L^2_z(\mathbb{G})$.

The Schur orthogonality relations \cite[Proposition~5.3.8~(iii)]{15} yield
\[
\langle \chi(\alpha),\chi(\beta) \rangle_{L^2}
   = \sum_{i=1}^{d_\alpha}\sum_{j=1}^{d_\beta} h\bigl((u_{ii}^\alpha)^* u_{jj}^\beta\bigr)
   = \delta_{\alpha,\beta}\sum_{i=1}^{d_\alpha} \frac{1}{d_\alpha}
   = \delta_{\alpha,\beta}.
\]
Hence $\chi$ sends the orthonormal basis $\{\delta_\alpha\}_{\alpha\in\operatorname{Irred}(\mathbb{G})}$ 
of $\ell^2(\operatorname{Irred}(\mathbb{G}))$ onto an orthonormal basis of 
$L^2_z(\mathbb{G})$, and therefore extends uniquely to a unitary operator
\[
\tilde{\chi}: \ell^2(\operatorname{Irred}(\mathbb{G})) \longrightarrow L^2_z(\mathbb{G}).
\]

\emph{Step 2: Intertwining of the representations.}  
Recall that $F(\mathbb{G})$ acts on $\ell^2(\operatorname{Irred}(\mathbb{G}))$ via 
the left regular representation $\Lambda_0$ of the fusion algebra, and that $\lambda$ 
denotes the left multiplication action of $\operatorname{Pol}(\mathbb{G})$ on 
$L^2(\mathbb{G})$. For any $\alpha,\beta\in\operatorname{Irred}(\mathbb{G})$ we have
\[
\tilde{\chi}\Lambda_0(\beta)(\delta_\alpha)
   = \tilde{\chi}(\beta\cdot\alpha)
   = \chi(\beta\cdot\alpha)
   = \lambda\bigl(\chi(\beta)\bigr)\chi(\alpha)
   = \lambda\bigl(\chi(\beta)\bigr)\tilde{\chi}(\delta_\alpha),
\]
where we used that $\chi$ is a unital $*$-isomorphism from $F(\mathbb{G})$ onto 
$\operatorname{Pol}_z(\mathbb{G})$ and that $\lambda$ is a faithful representation. 
By linearity this identity extends to all finitely supported functions, so for 
every $x\in F(\mathbb{G})$ we obtain
\[
\Lambda_0(x) = \tilde{\chi}^* \lambda(\chi(x)) \tilde{\chi}
\]
as operators on the dense subspace $c_c(\operatorname{Irred}(\mathbb{G}))$. 
Consequently $\Lambda_0(x)$ extends to a bounded operator on 
$\ell^2(\operatorname{Irred}(\mathbb{G}))$, and this extension satisfies 
$\Lambda(x)=\tilde{\chi}^*\,\lambda(\chi(x))\,\tilde{\chi}$.

\emph{Step 3: Intertwining of the Dirac operators.}  
For any $\alpha\in\operatorname{Irred}(\mathbb{G})$ we compute
\[
\tilde{\chi}\tilde{D}_\ell(\delta_\alpha) = \ell(\alpha)\,\tilde{\chi}(\delta_\alpha)
                                   = \ell(\alpha)\,\chi(\alpha)
                                   = D_\ell^z\,\tilde{\chi}(\delta_\alpha),
\]
where $\ell(\alpha)$ denotes the length of $\alpha$ (i.e., the eigenvalue of the 
Dirac operator on the basis vector $\delta_\alpha$). Thus $\tilde{\chi}$ maps the 
core $\operatorname{span}\{\delta_\alpha\}$ into $\operatorname{Dom}(D_\ell^z)$ 
and satisfies $\tilde{\chi}\tilde{D}_\ell = D_\ell^z\tilde{\chi}$ on this core. 
It follows immediately that $\tilde{\chi}(\operatorname{Dom}(\tilde{D}_\ell))
\subset\operatorname{Dom}(D_\ell^z)$ and the equality holds on the whole domain. 
Hence $\tilde{\chi}$ is a unitary equivalence between the two spectral triples.
\end{proof}

\begin{theorem}\label{Theorem5.5}
  Let $\mathbb{G}$ be a coamenable compact quantum group equipped with a proper length 
  function. If $\bigl(\operatorname{Pol}_z(\mathbb{G}), L_{z,\ell}^k\bigr)$ 
  is a compact quantum metric space and the  truncations $\bigl(P_\Lambda\operatorname{Pol}_z(\mathbb{G}), L_{z,\ell,\Lambda}^k\bigr)$ of 
  $\bigl(\operatorname{Pol}_z(\mathbb{G}), L_{z,\ell}^k\bigr)$ exhibit 
  Gromov--Hausdorff convergence, then 
  \[
  \operatorname{dist}_q\bigl(F(\mathbb{G}), L_{\ell}^k),\, 
  (P_\Lambda F(\mathbb{G}), L^k_{\ell,\Lambda})\bigr) \to 0 \qquad as \quad\Lambda\to \infty.
  \]
\end{theorem}

\begin{proof}
 By Proposition~\ref{prop unitary}, we know 
 $\bigl(\operatorname{Pol}_z(\mathbb{G}), L_{z,\ell}^k\bigr)$ is a compact 
 quantum metric space if and only if $(F(\mathbb{G}), L_{\ell}^k)$ is a compact 
 quantum metric space. Let 
 $\tilde{\chi}_\Lambda^* \colon P_\Lambda(L^2_z(\mathbb{G})) \to 
 P_\Lambda(\ell^2(\operatorname{Irred}(\mathbb{G})))$ be the unitary operator 
 obtained by restricting $\tilde{\chi}^*$ to the truncated subspaces; it is 
 therefore an isometry (in fact, unitary).

 By hypothesis, the  truncations $\bigl(P_\Lambda\operatorname{Pol}_z(\mathbb{G}), L_{z,\ell,\Lambda}^k\bigr)$ of 
 $\bigl(\operatorname{Pol}_z(\mathbb{G}), L_{z,\ell}^k\bigr)$ converge in quantum 
 Gromov--Hausdorff distance.

 From Proposition~\ref{prop unitary}, $\tilde{\chi}$ is a unitary equivalence 
 between the above spectral triples:
 \[
 \tilde{\chi}^{*}\bigl(\operatorname{Pol}_z(\mathbb{G}), L^2_z(\mathbb{G}), D_\ell^z\bigr)
 \tilde{\chi} = \bigl(F(\mathbb{G}), \ell^2(\operatorname{Irred}(\mathbb{G})), 
 \tilde{D}_\ell\bigr).
 \]
 Similarly, $\tilde{\chi}_\Lambda$ provides a unitary equivalence between the 
 truncated versions:
 \[
 \tilde{\chi}_\Lambda^{*}\bigl(P_\Lambda\operatorname{Pol}_z(\mathbb{G}) P_\Lambda,\,
 P_\Lambda L^2_z(\mathbb{G}),\, D_{z,\ell,\Lambda}\bigr) \tilde{\chi}_\Lambda 
 = \bigl(P_\Lambda F(\mathbb{G}) P_\Lambda,\, P_\Lambda\ell^2(\operatorname{Irred}(\mathbb{G})),\,
 \tilde{D}_{\ell,\Lambda}\bigr).
 \]

 Since unitary equivalence preserves the quantum Gromov--Hausdorff distance, 
 we obtain
 \[
 \operatorname{dist}_q\bigl(F(\mathbb{G}), L_{\ell}^k\bigr),\, 
 (P_\Lambda F(\mathbb{G}) , L^k_{\ell,\Lambda})\bigr) \to 0 \qquad as \quad\Lambda\to \infty.
 \]
\end{proof}

Since the spectral-truncated Lipschitz seminorm on the truncated space differs from $L_{\ell,\Lambda}^k$, the above theorem is not directly applicable to spectral truncations. However, by Theorem~\ref{compact quantum metric space} , Theorem~\ref{compact quantum metric space on fusion algebra} and  Corollary~C in \cite{16}, we can still construct compact quantum metric spaces whose Lip-norms are associated with proper length functions. Hence, we obtain examples whose spectral truncations exhibit Gromov--Hausdorff convergence.


\begin{example}\label{ex:spectral-truncation-GH}
   By Corollary~C of \cite{16}, we know that, when equipped with the length function $\ell$ arising from the standard fundamental corepresentation, the pairs
\[
    (C(SU(2)),L_{\ell}),\quad (C(O_{2}^{+}),L_{\ell}),\quad
    (C(SO(3)),L_{\ell}),\quad \text{and}\quad (C(S_{4}^{+}),L_{\ell})
\]
are compact quantum metric spaces.

Moreover, by Theorem~\ref{compact quantum metric space} and Theorem~\ref{compact quantum metric space on fusion algebra}, for every $N\in\mathbb{N}$ there exists $k\in\mathbb{N}$ such that $(SU(N),L_\ell^k)$ and $(SO(N),L_\ell^k)$ are compact quantum metric spaces.

Finally, by Theorem~\ref{theorem Gromov Hausdorff}, the spectral truncations of these quantum compact metric spaces, as well as those of their central function forms, converge in the Gromov--Hausdorff sense.
\end{example}

Now we consider another important class of examples: discrete quantum groups 
with the property of rapid decay. The study of this property for discrete 
quantum groups was initiated by Vergnioux \cite{20}. Building on this work, 
Bhowmick, Voigt, and Zacharias \cite{23} demonstrated that for quantum groups 
possessing the rapid decay property, one can naturally construct compact 
quantum metric spaces. In what follows, we briefly outline this construction.

Let $\hat{\mathbb{G}}$ be a discrete quantum group. We say that $\hat{\mathbb{G}}$ has \emph{property RD} 
with respect to a length function $\ell$ on $\hat{\mathbb{G}}$ if there exist constants 
$C, s > 0$ such that
\[
\| \mathcal{F}(f) \|_{\mathrm{op}} \leq C \| f \|_{2,s}
\]
for all $f \in C_c(\hat{\mathbb{G}})$. We say that $\hat{\mathbb{G}}$ has property RD if it has property RD 
with respect to some length function.

Let $k \in \mathbb{N}$ and define $T = C^{1/(2k)}$. For any 
$a \in \mathbb{C}(\hat{\mathbb{G}})$, set
\[
\delta(a) = [D_\ell, a],
\]
and define
\[
\delta_T(a) = [D_\ell, T a T], \qquad L_T^k(a) = \|\delta_T^k(a)\|.
\]
If $\hat{\mathbb{G}}$ is a discrete quantum group of rapid decay, by \cite{20} we 
know that $(C(\hat{\mathbb{G}}), L_T^k)$ is a compact quantum metric space.



\begin{theorem}\label{theorem GH on discrete}
Let $\hat{\mathbb{G}}$ be a discrete quantum group of rapid decay. Then 
\[
\operatorname{dist}_q\bigl((C(\hat{\mathbb{G}}), L_{T}^k),\, 
(P_\Lambda C(\hat{\mathbb{G}}) P_\Lambda,L_{C(\hat{\mathbb{G}})_\Lambda}^{\alpha^\tau,\beta^\tau})\bigr) \to 0 \qquad as \quad\Lambda\to \infty
\]
\end{theorem}

\begin{proof}
    The result follows by repeating the argument in the proof of 
    Theorem~\ref{theorem Gromov Hausdorff}.
\end{proof}

\begin{example}
    If $\hat{\mathbb{G}}$ is a discrete quantum group of rapid decay, by \cite{20} we know that $(C(\hat{\mathbb{G}}), L_T^k)$ is 
    a compact quantum metric space and $T$ commutes with $D$. So $L_T^k$ is 
    bi-invariant. By Theorem~\ref{theorem GH on discrete}, their spectral 
    truncation converges in the Gromov--Hausdorff distance.
\end{example}

























\bibliographystyle{amsplain}
\begin{thebibliography}{99}

\bibitem{1}Alain Connes, Walter D. van Suijlekom, Spectral truncations in noncommutative geometry and
operator systems, Commun. Math. Phys. 383 (3) (2021) 2021--2067.


\bibitem{2}   Guixiang Hong, Simeng Wang, Xumin Wang, Pointwise convergence of noncommutative Fourier
series, Mem. Am. Math. Soc. 302 (1520) (2024), v+86.

\bibitem{3} David Kyed, L2-Betti numbers of coamenable quantum groups, Münster J. Math. 1 (2008) 143--179.
\bibitem{4} Malte Leimbach, Walter D. van Suijlekom, Gromov-Hausdorff convergence of spectral truncations
for tori, Adv. Math. 439 (2024) 109496.
\bibitem{5} Walter D. van Suijlekom, Gromov-Hausdorff convergence of state spaces for spectral truncations,
J. Geom. Phys. 162 (2021) 104075.
\bibitem{6} Malte Leimbach, Convergence of Peter-Weyl truncations of compact quantum groups, \textit{J. Noncommut. Geom.} online first, doi:10.4171/jncg/634.
\bibitem{7}Yvann Gaudillot-Estrada, Walter D. van Suijlekom, Convergence of spectral truncations for compact
metric groups, Int. Math. Res. Not. (13) (2025) rnaf197.
\bibitem{8}Hanfeng Li, Compact quantum metric spaces and ergodic actions of compact quantum groups, J.
Funct. Anal. 256 (10) (2009) 3368--3408.
\bibitem{9}Marc A. Rieffel, Gromov-Hausdorff distance for quantum metric spaces, Mem. Am. Math. Soc.
168 (796) (2004) 1--65. Appendix 1 by Hanfeng Li.
\bibitem{10} Marc A. Rieffel, Matrix algebras converge to the sphere for quantum Gromov-Hausdorff distance,
Mem. Am. Math. Soc. 168 (796) (2004) 67--91.
\bibitem{11}Marc A. Rieffel, Convergence of Fourier truncations for compact quantum groups and finitely generated groups, J. Geom. Phys. 192 (2023) 104921.
\bibitem{12}Stanisław L. Woronowicz, Compact quantum groups, in: Symétries Quantiques, Les Houches, 1995,
North-Holland, Amsterdam, 1998, pp. 845--884.

\bibitem{13}D. Burago, Y. Burago, and S. Ivanov. A course in metric geometry, volume 33 of Graduate
Studies in Mathematics. American Mathematical Society, Providence, RI, 2001.

\bibitem{14}Fumio Hiai and Masaki Izumi. Amenability and strong amenability for fusion algebras with appliations to subfator theory. Internat. J. Math.,
9(6):669 722, 1998.

\bibitem{15}Marc A. Rieffel, Group C*-algebras as compact quantum metric spaces, Doc. Math. 7 (2002)
605--651.

\bibitem{16}Are Austad and David Kyed.Quantum metrics from length function on quantum groups,J. Funct. Anal.290(2026)111256

\bibitem{17}Teo Banica,Introduction to Quantum Groups,Springer,Cham,2022.

\bibitem{18}Narutaka Ozawa, Marc A. Rieffel, Hyperbolic group C*-algebras and free-product C*-algebras as compact quantum metric spaces, Can. J. Math. 57 (5) (2005) 1056--1079.

\bibitem{19}Simeng Wang, Lacunary Fourier series for compact quantum groups, Commun. Math. Phys. 349 (3)
(2017) 895--945.

\bibitem{20}Roland Vergnioux,The property of rapid decay for discrete quantum groups.
J. Operator Theory 57 (2007), no. 2, 303–324.

\bibitem{21}Cristina Antonescu and Erik Christensen. Metrics on group C*-algebras and a non commutative Arzel'a-Ascoli theorem. J. Funct. Anal., 214(2):247-259, 2004.

\bibitem{22}Marc A. Rieffel, Compact quantum metric spaces, in: Operator Algebras, Quantization, and Non commutative Geometry, in: Contemp. Math., vol. 365, Amer. Math. Soc., Providence, RI, 2004,
pp. 315--330.

\bibitem{23}Jyotishman Bhowmick, Christian Voigt, Joachim Zacharias, Compact quantum metric spaces from quantum groups of rapid decay, J. Noncommut. Geom. 9 (4) (2015) 1175--1200.

\bibitem{24}Alain Connes. Compact metric spaces, Fredholm modules, and hyperfiniteness. Ergodic Theory and Dynamical Systems, 9(2):207-220, 1989.

\bibitem{25}Michael Christ and Marc A Rieffel. Nilpotent group C*algebras as compact quantum metric spaces. Canadian Mathematical Bulletin, 60(1):77-94, 2017.

\bibitem{26}Thomas Timmermann, An Invitation to Quantum Groups and Duality, EMS Textbooks in Mathematics, European Mathematical Society (EMS), Zürich, 2008. From Hopf algebras to multiplicative unitaries and beyond.

\bibitem{27}Jens Kaad, External products of spectral metric spaces, Pure Appl. Funct. Anal. 9 (5) (2024)
1175--1211


\end{thebibliography}
\end{document}


%------------------------------------------------------------------------------
% End of journal.tex
%------------------------------------------------------------------------------
